\documentclass[11pt,a4paper]{article}
\usepackage[T1]{fontenc}
\usepackage[utf8]{inputenc}
\usepackage{lmodern}
\usepackage[margin=27mm,headheight=15pt]{geometry}
\usepackage{amsmath,amssymb,amsthm,mathtools}
\usepackage{microtype}
\usepackage{booktabs,longtable,array,enumitem}
\usepackage{xcolor}
\usepackage{fancyhdr}
\usepackage[colorlinks=true,linkcolor=blue!45!black,citecolor=blue!45!black,urlcolor=blue!55!black]{hyperref}
\usepackage{bookmark}
\hypersetup{pdftitle={Gap-Six Reconstruction, Complex Harmonic Determinants, and Stieltjes Boundary Collisions},
 pdfauthor={Research continuation for the ProveIt project},
 pdfsubject={Exact identities, proofs, and reproducible mathematical certificates}}
\newtheorem{theorem}{Theorem}[section]
\newtheorem{lemma}[theorem]{Lemma}
\newtheorem{proposition}[theorem]{Proposition}
\newtheorem{corollary}[theorem]{Corollary}
\theoremstyle{definition}
\newtheorem{definition}[theorem]{Definition}
\theoremstyle{remark}
\newtheorem{remark}[theorem]{Remark}
\numberwithin{equation}{section}
\setlist{itemsep=4pt,topsep=5pt}
\setlength{\emergencystretch}{2em}
\allowdisplaybreaks[2]
\pagestyle{fancy}
\fancyhf{}
\fancyhead[L]{\small ProveIt research continuation}
\fancyhead[R]{\small 12 October 2026}
\fancyfoot[C]{\thepage}
\renewcommand{\headrulewidth}{0.3pt}
\title{\LARGE\bfseries Gap-Six Reconstruction,\\[3pt]
 Complex Harmonic Determinants,\\[3pt]
 and Stieltjes Boundary Collisions\\[10pt]
 \large Exact identities and reproducible proofs}
\author{Research continuation for the ProveIt project}
\date{12 October 2026\\[3pt]\small Repository snapshot: 11 October 2026}
\begin{document}
\maketitle
\begin{abstract}
We resolve three questions left open in the inspected ProveIt
contributions and develop explicit identities from the resulting
structure. First, every unequal even harmonic-tail pair with gap
six has the optimal consecutive moment prefix. The proof makes
the existing Gamma-basis asymptotic effective above an explicit
threshold and certifies the entire finite complement by exact
modular arithmetic. Second, formal harmonic degree equals generic
pencil rank for every nonzero complex rational weight. A short
rational difference theorem for symmetric matrix primitives removes
the real-positivity hypothesis and allows varying pencil kernels.
Third, we determine the full power--logarithm expansion, including
the constant, for arbitrary products of argument derivatives of
generalized Stieltjes constants whose positive shifts approach zero
at unequal fixed rates.

The identity consequences include an all-order family of ordinary
convergent rational polygamma-tail sums; a three-by-three Hahn
formula for highest harmonic coefficients at arbitrary particle
count; an exact classification of adjacent-pole residue exceptions;
cutoff generators using only rational functions, logarithms, and
one classical dilogarithm; and an explicit positive three-particle
ensemble integral evaluated in ordinary zeta values. Proofs,
finite exact certificates, and numerical diagnostics are distinguished
throughout. The unrestricted all-pairs moment conjecture and the
manuscript's short Gaussian $S_6$ and revised $S_8$ candidates
remain open.
\end{abstract}
\clearpage
\begingroup
\small
\tableofcontents
\endgroup
\clearpage
% BEGIN sections/00_scope.tex
\section{Scope, source frontier, and principal results}
\label{scope:section}

This article continues the program in the two directories specified
by the project owner: the polylogarithm manuscript and the incoming
research contributions \cite{RepoManuscript,RepoIncoming}. All
repository assertions refer to the pinned revision
\begin{center}\small\ttfamily
c4036d60715916ae72dd18983782fdb182840b86.
\end{center}
The incoming index at that revision contains eleven archives.
We extracted them, reconciled their overlapping research-status
lists, and read the proof dependencies used below. This matters:
an open question in an earlier contribution can already be solved
in a later archive at the same revision.

In particular, the smaller-order-ten and smaller-order-twelve
moment families had already been completed, as had the complete
gap-two and gap-four families, several Stieltjes spectral faces,
and arbitrary mixed even-tail products
\cite{BilinearIncoming,TailKernelsIncoming,OrderedIncoming,LatticeIncoming}.
Those results are inputs or context, and are not claimed anew.
The useful remaining targets were the complete gap-six family,
complex rational harmonic degree, and unequal-rate boundary
collisions. The new theorems address these three targets directly.

\subsection{What changes at this frontier}

\begin{longtable}{@{}>{\raggedright\arraybackslash}p{0.23\textwidth}
 >{\raggedright\arraybackslash}p{0.40\textwidth}
 >{\raggedright\arraybackslash}p{0.29\textwidth}@{}}
\toprule
Question at the pinned revision & Result in this article & Precise scope\\
\midrule
\endfirsthead
\toprule
Question at the pinned revision & Result in this article & Precise scope\\
\midrule
\endhead
Optimal moment prefix for all even pairs with gap six &
Theorem~\ref{gs:main}: the first $q$ centered moments reconstruct
all moments for $p=q-6$, $q\ge8$ even. &
Full formal coefficient rank; an analytic infinite range plus an
exact finite certificate.\\[7pt]
Harmonic degree for complex rational weights &
Theorem~\ref{cr:degree}: degree equals generic pencil rank for
every nonzero complex rational weight. &
Rational dependence on cutoff; only finitely many admissible
cutoff exceptions. No positivity hypothesis.\\[7pt]
Explicit adjacent-pole determinants beyond small rank &
Theorems~\ref{hp:hahn-top} and \ref{hp:all-r-generator}: a
three-by-three Hahn formula and a classical-polylogarithm
generator for every particle count. &
Specified two-pole family, with every endpoint residue included.\\[7pt]
Unequal-rate Stieltjes boundary collisions &
Theorems~\ref{urc:constant-theorem} and \ref{urc:full-theorem}:
all index and argument-derivative orders, with a finite contact
formula for the constant. &
Fixed positive rates, and locally uniform joint spectral
subtraction.\\
\bottomrule
\end{longtable}

Several identity consequences have a particularly short form.
For example, with $H_M=\sum_{k=1}^M1/k$, the positive ensemble
in Section~\ref{hp:section} satisfies
\[\begin{aligned}
 &\sum_{0\le n_1<n_2<n_3<M}
 \frac{\prod_{i<j}(n_j-n_i)^2}{\prod_i(n_i+1)(n_i+2)}\\
 &\qquad=
 -\frac{M(M+1)(M+2)}3H_M^2+
 \frac{4M^2(M+2)}3H_M+
 \frac{M^3(M^2-18M-31)}{12(M+1)}.
\end{aligned}\]
The boundary calculation produces the equally compact ordered
identity
\[
 \operatorname{CT}_{\epsilon\downarrow0}
 \int_0^1\psi'(x+\lambda\epsilon)\psi(x+\mu\epsilon)\,dx
 =\zeta(2)\left(\frac\mu\lambda+\log\frac\mu\lambda\right),
 \qquad\lambda,\mu>0.
\]
Here $\operatorname{CT}$ is explicitly the coefficient of
$\epsilon^0(\log\epsilon)^0$, not the value of a divergent
unsubtracted integral. The article supplies the complete
subtractions as well as this constant.

\subsection{The three mechanisms}

The first mechanism turns a limiting determinant into an exact
all-index result. The incoming Gamma basis has three exponential
scales for gap six. The middle scale loses its leading imaginary
phase, so a one-term Laplace estimate cannot prove the theorem.
Retaining the next term produces a normalized reciprocal-square
Vandermonde. A quantitative error comparison proves nonvanishing
for every $b\ge2048$, and exact recurrence arithmetic covers
$4\le b\le2047$. Estimates are used only to finish an identity and
reconstruction problem; no inequality is presented as the main
research output.

The second mechanism is algebraic. If a symmetric rational matrix
$G(z)$ has a rank-one difference
\[
 G(z)-G(z-1)=w(z)p(z)p(z)^{\mathsf T}
\]
with $w\ne0$ and independent rational components of $p$, then
$G$ cannot be generically singular. Otherwise its rank stabilizes
at large integer cutoffs, equal ranks force its image to remain
fixed, and the update vectors lie in a fixed proper subspace.
That would make a nonzero constant combination of the components
of $p$ vanish identically. This elementary argument works over
every field of characteristic zero, including complex fields with
isotropic vectors.

The third mechanism is local at the singular endpoint. The
Hurwitz recurrence splits each Stieltjes generating function into
one explicit singular power and an analytic function at one.
Only a critical Taylor coefficient of the analytic factor can
contribute to the rate-dependent constant. The singular--singular
term carries a beta or polylogarithmic kernel but contributes
strictly negative powers of the collision parameter. The resulting
constant is therefore much smaller than the full two-shift
integral.

\subsection{Conventions and logical status}

We use $B_1=-1/2$, rising factorials $(a)_k$, and falling
factorials $a^{\underline k}$. Generalized harmonic numbers are
$H_N^{(s)}(a)=\sum_{n=0}^N(n+a)^{-s}$, while the unshifted
notation $H_M$ runs from $1$ to $M$. The polygamma convention is
$\psi^{(r)}=d^{r+1}\log\Gamma/dz^{r+1}$, so
$\psi^{(r)}(z)=(-1)^{r+1}r!\zeta(r+1,z)$ for $r\ge1$.
The generalized Stieltjes constants are fixed by
\[
 \zeta(1+u,z)=\frac1u+
       \sum_{m\ge0}\frac{(-1)^m}{m!}\gamma_m(z)u^m.
\]
Unless otherwise stated, argument derivatives are indicated by
parenthesized superscripts; spectral derivatives are written
explicitly.

Formal harmonic coordinates and formal zeta labels are independent
symbols in their specified coefficient spaces. A rank theorem for
that space is not a theorem of arithmetic independence for its
evaluated constants. Similarly a cutoff generating function is
first an identity on $|x|<1$; a boundary value or an antiderivative
at one is given its own convergence argument.

The article contains four kinds of evidence. Analytic and
algebraic proofs establish the general statements. Exact modular
nonvanishing is part of the gap-six proof on a finite, explicitly
bounded interval. Other exact finite checks verify symbolic
formulas and normalizations, without being substituted for general
proofs. High-precision quadratures are diagnostics, with no claim
that floating-point residuals are certified enclosures.

The methods of Christoffel transformations, Hahn polynomials,
rational discrete residues, Euler sums, and local finite-part
subtraction are established. They are credited to primary
references where used \cite{Krattenthaler2023,DLMF,ChenSinger2012,
Xu2017,XuWang2019,ShpotParis2016}. The claimed advances are the
precise theorems and specializations relative to the inspected
ProveIt frontier. The literature review supports that distinction
but is not an exhaustive historical-priority search. Independent
internal proof audits are included; external peer review and
proof-assistant formalization remain separate tasks.


% END sections/00_scope.tex

% BEGIN sections/01_gap_six.tex
% Parent-ready proof fragment. All local labels use gs:.
% Mathematical inputs inspected at c4036d60715916ae72dd18983782fdb182840b86:
% Ordered_Boundary_Identities..., sections/08_gamma_basis.tex;
% harmonic-tail-gamma, sections/02_gap_four.tex.
% The all-gap-six theorem and explicit threshold below are the new step.
\section{The complete gap-six harmonic-tail reconstruction theorem}
\label{gs:section}

The fixed-gap Gamma basis of the incoming \emph{Ordered Boundary
Identities and Gamma Kernels} report~\cite{OrderedIncoming} proves eventual reconstruction
on every fixed-gap ray but does not specify a threshold.  We make its
argument effective when the gap is six and complete the remaining
finite interval by exact modular arithmetic.  This proves the entire
previously unresolved gap-six family.  The finite part is a
computation-assisted proof: every determinant is certified by a
nonzero residue in a finite field, with a short independently runnable
certificate.  No floating-point approximation enters that certificate.

\subsection{Statement and the exact rational obstruction}

Let $b\ge4$ be an integer, $p=2b-6$, $q=2b$, and
\[
 A_p(n)=\zeta(p)-H_n^{(p)}=\zeta(p,n+1),\qquad x_n=n+\tfrac12.
\]
Let $D_r=D_{p,q}(-2r)$, where
\[
 D_{p,q}(s)=\sum_{n\ge0}A_p(n)A_q(n)x_n^{-s}
\]
is initially defined in its convergence half-plane and subsequently
continued meromorphically.  The centered Hurwitz expansion shows that
its possible poles occur at odd integers, so all displayed $D_r$ are
regular.  The inherited moment blocks satisfy the inverse transforms
of~\cite[Section~8]{OrderedIncoming}:
\begin{align}
 D_r&=\sum_{j=0}^r\frac{\binom{2r}{2j}}{2j+1}
           B_{2r-2j}(1/2)C_j,\label{gs:forward}\\
 C_j&=\sum_{r=0}^j\binom{2j+1}{2r}2^{-2(j-r)}D_r.
 \label{gs:backward}
\end{align}
For $j\ge b$, write
\begin{equation}\label{gs:late}
 C_j=\sum_{k=0}^{b-1}a_{jk}Z_k,
 \qquad Z_k=\zeta(2k),\qquad Z_0=-\tfrac12.
\end{equation}
With $d=2j+1$ and $e=d-p-q+2k+1$, these rational entries are zero
when $e<0$, and otherwise the exact late-coordinate formula of
\cite[Section~8]{OrderedIncoming} is
\begin{equation}\label{gs:entry}
 a_{jk}=B_e\left\{
  \frac{\binom{d-q+1}{p-2k}}{d-q+1}
 +\frac{\binom{d-p+1}{q-2k}}{d-p+1}\right\}.
\end{equation}
Every binomial summand whose lower index is not between $2$ and its
upper index is omitted.  In particular a lower index of zero is
omitted.  Put
\[
 M_b=(a_{b+r,k})_{0\le r,k<b}.
\]
All ranks below are ranks of rational coefficient matrices on the
specified formal labels.  They do not assert arithmetic independence
of the actual zeta values.

\begin{theorem}[Every gap-six pair has the optimal prefix]
\label{gs:main}
For every integer $b\ge4$, the matrix $M_b$ is nonsingular.  Therefore
\begin{align}\label{gs:span}
 \operatorname{span}_{\mathbb Q}\{D_0,D_1,\ldots\}
 &=\operatorname{span}_{\mathbb Q}\{D_0,\ldots,D_{2b-1}\}\\
 &=\operatorname{span}_{\mathbb Q}
 \{C_0,\ldots,C_{b-1},Z_0,\ldots,Z_{b-1}\}.\notag
\end{align}
The first $q=2b$ moments have formal rank $q$, and their coefficient
determinant in the displayed coordinate order is
$\det M_b/(4b-1)!!$.  This number of consecutive moments is optimal
for rank $q$ on those $q$ formal labels.
\end{theorem}

Define, for $n\ge1$,
\[
 \alpha_n=\frac{2|B_{2n}|}{(2n)!},\qquad
 A_u=\begin{cases}0,&1\le u\le3,\\
                   (2u)^{\underline6},&u\ge4,
       \end{cases}
\]
and, for $s\ge4$,
\[
 B_s=(2s-1)^{\underline6},\qquad d_s=2(A_s+B_s)>0.
\]
Here the underline denotes a falling factorial.  The $B_s$ in this
line are recurrence factors, distinguished from the Bernoulli numbers
by the displayed definition.  Define a rational array by
\begin{align}
 R_{s,j}&=\delta_{sj}&&(1\le s,j\le3),\notag\\
 R_{s,j}&=\frac1{d_s}\sum_{u=1}^{s-1}
             \alpha_{s-u}(B_s+A_u)R_{u,j}
           &&(s\ge4,\ 1\le j\le3),\label{gs:recurrence}
\end{align}
and set
\[
 F_b=(R_{b+i,j})_{1\le i,j\le3},\qquad f_b=\det F_b.
\]
The inherited complementary-minor identity
\cite[Section~8, fixed-gap boundary reduction]{OrderedIncoming} is
\begin{equation}\label{gs:bridge}
 \det M_b=\frac{(-1)^{b(b-1)/2}}{4^bb!}
        \left(\prod_{s=4}^{b+3}d_s\right)f_b.
\end{equation}
Thus the theorem reduces exactly to $f_b\ne0$ for every $b\ge4$.

For clarity, the finite differential mechanism behind this reduction
is the operator
\[
 \mathcal T R=(R^{(6)}-R^{(6)}(0))\cot(z/2)
                       +(R\cot(z/2))^{(6)}.
\]
The coefficient of $z^{2s-7}$ in $\mathcal T R$, for
$R(z)=\sum_{u\ge1}r_uz^{2u}$, is
\[
 d_sr_s-\sum_{u<s}\alpha_{s-u}(B_s+A_u)r_u.
\]
The first three $r_u$ are free and every later one is determined.
Truncating these three null solutions at degree $2b$ and applying
Jacobi's complementary-minor identity gives \eqref{gs:bridge}; the
row and reversed-column factorial scales contribute $4^bb!$ and
$(-1)^{b(b-1)/2}$, respectively.  This is the same normalization as
\eqref{gs:entry}, including its omitted zero lower index.

\subsection{The Gamma representation needed for an effective estimate}

Put
\begin{align*}
 W(x)&=\left|\Gamma\!\left(\tfrac14+\tfrac{ix}{2}\right)\right|^{12},\\
 P_1(x)&=\sinh(2\pi x),\qquad
 P_2(x)=\sinh(\pi x),\qquad
 P_3(x)=\tanh(\pi x),
\end{align*}
and define
\begin{equation}\label{gs:t}
 t_{n,j}=\frac4{(2n)!}\operatorname{Im}
           \int_0^\infty(x+i/2)^{2n}P_j(x)W(x)\,dx.
\end{equation}
All these integrals converge absolutely.  The inherited Gamma-basis
result~\cite[Section~8, Gamma basis and positive normalization]{OrderedIncoming}
says that the functions
\[
 G_j(z)=4\sin(z/2)\int_0^\infty\sinh(xz)P_j(x)W(x)\,dx
\]
satisfy $\mathcal T G_j=0$, and their first three coefficient rows
form a matrix $T=(t_{n,j})_{1\le n,j\le3}$ with $\det T<0$.
Consequently
\begin{equation}\label{gs:Gamma-bridge}
 (t_{b+i,j})_{i,j=1}^3=F_bT,\qquad
 f_b=\frac{\det(t_{b+i,j})_{i,j=1}^3}{\det T}.
\end{equation}

These inputs have short verification mechanisms.  The null equation
follows from the Gamma recurrence after shifting two Fourier-integral
contours by $\pm i$.  The shifted common Gamma factor has its nearest
poles at $\pm3i/2$; only the tanh multiplier contributes the two
simple poles at $\pm i/2$.  Their residue terms provide exactly the
subtracted $G_3^{(6)}(0)$, while the sinh multipliers contribute zero.
For the initial determinant, $\operatorname{Im}(x+i/2)^{2n}$ has
leading term $nx^{2n-1}$.  Continuous Cauchy--Binet therefore reduces
its sign to the product of the two alternants
\[
 \det(x_r^{2n-1})_{n,r=1}^3,
 \qquad \det(P_j(x_r))_{r,j=1}^3.
\]
After extracting $\prod_r x_r\tanh(\pi x_r)$, the first is a
Vandermonde in $x_r^2$ and the second is minus twice a Vandermonde
in $\cosh(\pi x_r)$.  Both variable maps are strictly increasing
on $(0,\infty)$; integration against $\prod_rW(x_r)$ is therefore
strictly negative.  This proves $\det T<0$ without a numerical
normalization constant.

\subsection{An explicit density bound}

Let
\begin{equation}\label{gs:scales}
 \lambda_j=j\pi,\qquad A=8(2\pi)^6,
 \qquad a_1=a_2=A/2,\quad a_3=A.
\end{equation}

\begin{lemma}[Uniform Gamma density estimate]\label{gs:density}
For every $x\ge2$ and $1\le j\le3$,
\begin{equation}\label{gs:density-bound}
 P_j(x)W(x)=a_jx^{-3}e^{-\lambda_jx}(1+\delta_j(x)),
 \qquad |\delta_j(x)|\le\frac6{x^2}.
\end{equation}
Moreover $a_j>1$, and on $0\le x\le2$,
\begin{equation}\label{gs:small-density}
                  0\le P_j(x)W(x)/a_j<2^{54}.
\end{equation}
\end{lemma}

\begin{proof}
We use the classical complex Stirling remainder in DLMF
\S5.11(ii).  For $z=1/4+ix/2$, $|\arg z|<\pi/2$, so retaining
$1/(12z)$ in Stirling's logarithmic expansion gives
\[
 \log\Gamma(z)=(z-\tfrac12)\log z-z+\tfrac12\log(2\pi)
                         +\frac1{12z}+\rho(z),
 \qquad |\rho(z)|\le\frac1{90|z|^3}.
\]
Indeed the first omitted term has magnitude $1/(360|z|^3)$,
and the secant factor is at most $4$.
Set $t=1/(2x)$ and
\[
 L(x)=\log\frac{W(x)}{Ax^{-3}e^{-3\pi x}}.
\]
Taking twelve times the real part and simplifying the elementary
logarithm yields
\[
 L(x)=-\frac32\log(1+t^2)
       +3\left(\frac{\arctan t}{t}-1\right)
       +\frac{4t^2}{1+t^2}+12\operatorname{Re}\rho(z).
\]
The inequalities $\log(1+t^2)\le t^2$ and
$0\le1-\arctan(t)/t\le t^2/3$ give
\[
 |L(x)|\le\frac{13}{8x^2}+\frac{16}{15x^3}
        \le\frac{259}{120x^2}<\frac9{4x^2}\qquad(x\ge2).
\]
Since $e^{9/16}<2$,
\[
 \left|\frac{W(x)}{Ax^{-3}e^{-3\pi x}}-1\right|
        \le\frac9{2x^2}.
\]
The elementary multipliers differ from their leading exponentials by
at most $2e^{-2\pi x}$.  More precisely their normalized factors are
$1-e^{-4\pi x}$, $1-e^{-2\pi x}$, and
$(1-e^{-2\pi x})/(1+e^{-2\pi x})$, respectively, and each lies
between zero and one.  Since $\pi>3$ and
$e^{6x}\ge(6x)^2/2$,
\[
                    2e^{-2\pi x}\le\frac1{9x^2}.
\]
Combining the two errors proves \eqref{gs:density-bound}, with
room to spare.

Euler's Gamma integral gives
$|\Gamma(1/4+ix/2)|\le\Gamma(1/4)<5$, because the contributions
from $(0,1)$ and $(1,\infty)$ are less than $4$ and $1$.
For $0\le x\le2$ all three multipliers are less than $e^{16}$.
Using $e<3$, $5^{12}<2^{28}$, $3^{16}<2^{26}$, and $a_j>1$
proves \eqref{gs:small-density}.
\end{proof}

\subsection{A completely explicit shifted Laplace estimate}

For an even integer $N\ge4096$ and $\lambda\in\{\pi,2\pi,3\pi\}$,
write
\[
 S_N(\lambda)=\frac{\Gamma(N-2)}{\lambda^{N-2}},\quad
 S_N^{(1)}(\lambda)=\frac{\Gamma(N-3)}{\lambda^{N-3}},\quad
 S_N^{(2)}(\lambda)=\frac{\Gamma(N-4)}{\lambda^{N-4}}.
\]

\begin{lemma}[Effective two-term phase estimate]\label{gs:Laplace}
For $1\le j\le3$ and every even $N\ge4096$,
\begin{equation}\label{gs:Laplace-bound}
 \left|\frac1{a_j}\int_0^\infty(x+i/2)^NP_j(x)W(x)\,dx
 -e^{i\lambda_j/2}
     \left(S_N(\lambda_j)+\frac{3i}{2}S_N^{(1)}(\lambda_j)\right)
 \right|\le16S_N^{(2)}(\lambda_j).
\end{equation}
\end{lemma}

\begin{proof}
The contribution of $[0,2]$ is at most $2^{55}3^N$ by
\eqref{gs:small-density}.  On $[2,\infty)$ substitute the model
$x^{-3}e^{-\lambda x}$ and use \eqref{gs:density-bound}.
The density error is at most
\[
 6\int_2^\infty|x+i/2|^Nx^{-5}e^{-\lambda x}\,dx.
\]
On $x\ge\sqrt N$,
\[
 |x+i/2|^N\le x^N\exp\left(\frac{N}{8x^2}\right)<2x^N,
\]
so this part is at most $12S_N^{(2)}(\lambda)$.  On
$2\le x\le\sqrt N$, its bound is
\[
                       \sqrt N(N+1)^{N/2}.
\]
Here we used $6x^{-5}\le6/32<1$ and the length at most $\sqrt N$.

For the model integral set $y=x+i/2$ and shift the ray to the
positive real axis beginning at $2$.  The integrand
\[
       e^{i\lambda/2}y^N(y-i/2)^{-3}e^{-\lambda y}
\]
has no pole in this strip.  The short vertical segment at $2$
is bounded by $3^N/16$, and the segment at infinity vanishes.
For real $y\ge2$, the binomial series gives
\[
 \left|(1-i/(2y))^{-3}-1-\frac{3i}{2y}\right|
                  \le\frac5{2y^2}.
\]
To check the constant, write $u=1/(2y)\le1/4$; then
$\sum_{k\ge2}\binom{k+2}{2}u^k\le10u^2$ since its quotient
by $u^2$ is increasing and is less than $10$ at $u=1/4$.
Thus this model error is at most $(5/2)S_N^{(2)}(\lambda)$.
Extending the first two retained integrals from $[2,\infty)$
to $(0,\infty)$ contributes less than $2^N$.

It remains to absorb the compact and lower-range contributions.
Their sum is bounded by
\[
 R_N=2^{56}3^N+\sqrt N(N+1)^{N/2}.
\]
Since $\lambda<10$ and $(N-5)!\ge((N-5)/3)^{N-5}$,
\[
 S_N^{(2)}(\lambda)
 \ge\frac{(N-5)^{N-5}}{3^{N-5}10^{N-4}}.
\]
For every $N\ge4096$,
\[
 \frac{30\sqrt{N+1}}{N-5}\le\frac12,
 \qquad \frac{90}{N-5}\le\frac1{32},
 \qquad N^6\le2^{N-2}.
\]
The first follows by squaring and expanding a quadratic; the second
is immediate; the third holds at $N=4096$ and persists by induction
because $(1+1/N)^6<2$.  Therefore
\[
 \frac{\sqrt N(N+1)^{N/2}}{S_N^{(2)}}
       \le N^6 2^{-N}\le\frac14,
 \qquad
 \frac{2^{56}3^N}{S_N^{(2)}}
       \le2^{56}N^5 2^{-5N}<\frac14.
\]
Consequently $R_N<S_N^{(2)}$.  Adding the bounds gives
$12+5/2+1<16$, proving \eqref{gs:Laplace-bound}.
\end{proof}

The retained second term is essential.  At $j=2$ the leading
phase $e^{i\lambda_j/2}=-1$ has zero imaginary part, so discarding
that second term would give no nonzero leading coefficient for the
middle column.

\subsection{A determinant comparison valid for every large index}

Taking imaginary parts of \eqref{gs:Laplace-bound} gives, with
$N=2n\ge4096$,
\begin{align}
 t_{n,1}&=\frac{4a_1}{N!}S_N(\lambda_1)(1+\epsilon_{n,1}),\notag\\
 t_{n,3}&=-\frac{4a_3}{N!}S_N(\lambda_3)(1+\epsilon_{n,3}),
 &|\epsilon_{n,j}|&\le\frac{1600}{(N-4)^2}\quad(j=1,3),
 \label{gs:odd-errors}\\
 t_{n,2}&=-\frac{6a_2}{N!}S_N^{(1)}(\lambda_2)(1+\epsilon_{n,2}),
 &|\epsilon_{n,2}|&\le\frac{320}{3(N-4)}.
 \label{gs:even-error}
\end{align}
The estimates use only $\lambda_j<10$ and the exact Gamma ratios.

Let $b\ge2048$ and $n=b+i$ for $1\le i\le3$.
After extracting positive row and column factors and the column
signs $(+,-,-)$, the coefficient determinant is the determinant of
an array $E=(E_{ij})$ satisfying
\begin{equation}\label{gs:relative-matrix}
 E_{ij}=j^{-2i}(1+\eta_{ij}),\qquad |\eta_{ij}|<\frac1{25}.
\end{equation}
For precision, extract the positive row factors
$4\Gamma(2b+2i-2)\pi^{-2i}/\Gamma(2b+2i+1)$ and positive column
factors $a_j\lambda_j^{-2b+2}$ for $j=1,3$.
For $j=2$, multiply this column factor additionally by
$3\lambda_2/(2(2b-1))$.
The middle-column row correction is
\[
 \rho_i=\frac{2b-1}{2b+2i-3},\qquad
 0<\rho_i\le1,\qquad1-\rho_i\le\frac4{2b-1}.
\]
Thus its relative error is bounded by
\[
 \frac4{2b-1}+\frac{320}{3(2b-2)}<\frac1{25}\qquad(b\ge2048),
\]
while the outer-column bound follows at once from
\eqref{gs:odd-errors}.  The product of the three extracted signs is
positive.

For the exact limiting array $V_{ij}=j^{-2i}$ one computes
\begin{equation}\label{gs:V}
 \det V=-\frac5{1944},\qquad
 \operatorname{perm}V=\frac{289}{23328}.
\end{equation}
The permanent here is the sum of the six unsigned permutation
products.  By expanding both determinants,
\[
 |\det E-\det V|
 \le\operatorname{perm}V\left\{(26/25)^3-1\right\}
 =\frac{289}{23328}\frac{1951}{15625}<\frac5{1944}.
\]
The exact positive margin is $373661/364500000$.
Therefore $\det E<0$, and hence
\begin{equation}\label{gs:large-positive}
                         f_b>0\qquad(b\ge2048)
\end{equation}
by \eqref{gs:Gamma-bridge} and $\det T<0$.
This is an effective all-index conclusion, rather than a numerical
comparison at finitely many large indices.

\subsection{Exact finite completion}

The remaining range is $4\le b\le2047$.
The following two elementary recurrences give a finite certificate
without computing huge rational Bernoulli numbers.  From
\[
 \frac{x}{2}\cot\frac{x}{2}=1-\frac12\sum_{n\ge1}\alpha_nx^{2n}
\]
and the differential equation of the cotangent, one obtains
\begin{equation}\label{gs:alpha-recurrence}
 \alpha_1=\frac16,\qquad
 \alpha_n=\frac1{2(2n+1)}
       \sum_{k=1}^{n-1}\alpha_k\alpha_{n-k}\quad(n\ge2).
\end{equation}
Combine this recurrence with \eqref{gs:recurrence} modulo
\[
                         \ell=2147483647=2^{31}-1.
\]
The supplied script proves primality by trial division through
$\lfloor\sqrt\ell\rfloor$.  Every denominator is invertible in
this field: for the needed indices $n\le2049$ the factors
$2(2n+1)$ are smaller than $\ell$, and
\[
 d_s=2\frac{(2s)!}{(2s-6)!}\frac{2s-3}{s}
\]
for $4\le s\le2050$ has no factor divisible by $\ell$.
Induction therefore identifies every computed residue with the
reduction of its exact rational coefficient.

\begin{proposition}[Finite rational nonvanishing certificate]
\label{gs:finite}
For each integer $4\le b\le2047$, the determinant $f_b$ has a
nonzero reduction modulo $2147483647$.  Hence $f_b\ne0$ over
$\mathbb Q$ throughout this finite interval.
\end{proposition}

\begin{proof}[Exact computational proof]
Run \texttt{verify\_gap\_six.py}.  It constructs all needed
$\alpha_n$ by \eqref{gs:alpha-recurrence}, constructs the three
fundamental columns by \eqref{gs:recurrence}, and expands each
three-by-three determinant by its six signed products.  All 2,044
residues are nonzero; the complete indexed list is in
\texttt{gap\_six\_certificate.json}.  Its canonical compact JSON
residue list has SHA-256
\begin{center}\small\ttfamily
90ec1b66f919cd00404dd92f7d1f9483ceb9\\
d53bfa5893735db62c797113fea6
\end{center}
The hash identifies the record, while the recurrences and the nonzero
residues constitute the mathematical certificate.

As independent normalization checks, all 2,049 cotangent coefficients
agree with formal division of cosine by the normalized sine series,
and the bridge \eqref{gs:bridge} agrees with direct exact determinants
of the original Bernoulli matrices for $4\le b\le20$.  The latter
comparisons check the omitted-index convention and every row/column
factor against \eqref{gs:entry}.  They are additional checks; the
finite-field recurrence itself proves the entire finite interval.
\end{proof}

Combining \eqref{gs:large-positive}, Proposition~\ref{gs:finite}, and
\eqref{gs:bridge} proves Theorem~\ref{gs:main}.  The general
unrestricted unequal-even-order conjecture remains open beyond the
families now covered.  In particular, proving all gaps of six does
not by itself prove the gap-eight family.

\subsection{An all-order family of rational polygamma-tail evaluations}

The rank theorem turns the source block calculus into an exact
ordinary convergent identity for every gap-six pair.  Let
$(\lambda_b,\ldots,\lambda_{2b-1})$ be the first row of $M_b^{-1}$,
and put
\begin{align}
 K_b(y)&=\sum_{j=b}^{2b-1}\lambda_jy^{2j+1},\notag\\
 P_b(x)&=K_b(x+\tfrac12)-K_b(x-\tfrac12)
                  =\sum_{r=0}^{q-1}w_rx^{2r}.\label{gs:polynomial}
\end{align}
All coefficients are explicitly specified rational numbers; the
inverse exists for every $b\ge4$ by Theorem~\ref{gs:main}.
For $h=0,1,2,3$ define
\[
 c_{v,0}=\frac1{v-1},\qquad
 c_{v,h}=\frac{B_{2h}(1/2)}{(2h)!}(v)_{2h-1}\quad(h\ge1),
 \qquad u_h=\sum_{e=0}^hc_{p,e}c_{q,h-e}.
\]
Here $(v)_r$ denotes the rising factorial.  Define the even sextic
subtraction polynomial
\begin{equation}\label{gs:subtraction}
 S_b(x)=\sum_{h=0}^3\left(\sum_{e=0}^hu_e\,w_{q-1-h+e}\right)
                                      x^{6-2h}.
\end{equation}

\begin{corollary}[Rational evaluations for every even gap of six]
\label{gs:ordinary}
For $b\ge4$, $p=2b-6$, $q=2b$, and $x_n=n+1/2$,
\begin{equation}\label{gs:ordinary-identity}
 \boxed{\displaystyle
 \sum_{n=0}^{\infty}\left\{
 \frac{P_b(x_n)}{(p-1)!(q-1)!}
       \psi^{(p-1)}(n+1)\psi^{(q-1)}(n+1)-S_b(x_n)
                 \right\}=-\frac12.}
\end{equation}
This sum converges absolutely.  Its polynomial $P_b$ has exact degree
$2q-2$, with leading coefficient
\begin{equation}\label{gs:leading-coefficient}
 w_{q-1}=\frac{4b(4b-1)f_{b-1}}{d_{b+3}f_b}\ne0,
\end{equation}
where the auxiliary initial determinant is
$f_3=-1/638788869734400000$.
\end{corollary}

\begin{proof}
The first inverse row gives
$\sum_{j=b}^{2b-1}\lambda_jC_j=Z_0=-1/2$.
The coefficients of the centered difference in
\eqref{gs:polynomial} are exactly
\[
 w_r=\sum_{j=\max\{b,r\}}^{2b-1}
       \lambda_j\binom{2j+1}{2r}2^{-2(j-r)},
\]
so \eqref{gs:backward} gives $\sum_rw_rD_r=-1/2$.
The centered tail expansion gives
\[
 A_p(n)A_q(n)=\sum_{h=0}^3u_hx_n^{2-p-q-2h}
                       +O(x_n^{-p-q-6}).
\]
Since $\deg P_b\le2q-2$ and $q-p=6$, its nonnegative powers
are exactly \eqref{gs:subtraction}; the remaining summand is
$O(x_n^{-2})$.  The continued sums of $1,x_n^2,x_n^4,x_n^6$
are zero because $\zeta(-2h,1/2)=0$.  Thus the ordinary
absolutely convergent sum equals the displayed moment combination.
Finally $A_p(n)=\psi^{(p-1)}(n+1)/(p-1)!$ for even $p$.

For \eqref{gs:leading-coefficient}, suppose first that $b\ge5$.
The cofactor defining $(M_b^{-1})_{0,b-1}$ deletes the final row
and the first column of $M_b$.  Reindexing $j'=j-1$, $k'=k-1$,
$p'=p-2$, and $q'=q-2$ in \eqref{gs:entry} shows that this
minor is exactly $M_{b-1}$: both binomial pairs, the Bernoulli
index, and the omitted-index convention are unchanged.  Hence
\[
 \lambda_{2b-1}=(-1)^{b-1}\frac{\det M_{b-1}}{\det M_b}
              =\frac{4bf_{b-1}}{d_{b+3}f_b}
\]
by \eqref{gs:bridge}.  The centered difference multiplies its
highest coefficient by $4b-1$, proving the formula and
nonvanishing for $b\ge5$.  At $b=4$ a direct exact calculation
gives $\lambda_7=84/83$, $d_7=6795360$,
$f_4=-83/22789201752651350016000000$, and the stated value of
$f_3$; these verify the same formula without invoking a
zero-order tail.  This completes the proof.
\end{proof}

\paragraph{Further research.}
The same Gamma basis supplies $m$ distinct exponential scales for
any fixed half-gap $m$.  Making all of its phase estimates effective
would give a numerical finite-completion threshold for each fixed
gap.  The next concrete target is $m=4$, where the two even modes
require their first nonzero phase corrections and the determinant
comparison involves a four-by-four reciprocal-square Vandermonde.
A different question is whether a direct recurrence invariant can
prove gap-six nonvanishing without the finite completion; the initial
sign transition shows why a fixed-sign assertion from the first
admissible index is false.

% END sections/01_gap_six.tex

% BEGIN sections/02_complex_rank.tex
% Self-contained research fragment.  Requires amsmath, amssymb, amsthm.
% All labels use cr:.  References: BilinearIncoming, ChenSinger2012.
\section{Complex rational weights: a complete degree--rank theorem}
\label{cr:section}

The recent \emph{Bilinear Tails and Harmonic Determinants} contribution
proves that the formal harmonic degree of a moment determinant equals
its generic harmonic-pencil rank for every nonzero real rational weight.
It explicitly leaves the corresponding assertion for complex rational
weights open \cite[Section~5, Corollary \texttt{deg:rational};
Section~10, ``Complex weights and sharp cutoff exceptions'']{BilinearIncoming}.
We prove the complex assertion.  The new argument uses rationality in
the cutoff and the rigidity of symmetric rank-one updates.  It uses
ordinary transposition throughout; no Hermitian positivity is assumed.
The stronger numerical bound on the number of exceptions for eventually
one-sign real weights remains a separate advantage of the source theorem.

\subsection{A rational difference principle}

All fields in the next statements have characteristic zero.  The
indeterminate is $z$, and linear independence of functions means
independence over the constant field, unless another field is displayed.

\begin{lemma}[Equal ranks force equal images]\label{cr:rankone}
Let $A$ be a symmetric $q\times q$ matrix over a field $F$, let
$c\in F\setminus\{0\}$, and let $u\in F^q$.
If $u\notin\operatorname{im}A$, then
\[
 \operatorname{rank}(A+cuu^{\mathsf T})=\operatorname{rank}A+1.
\]
Consequently, if $A$ and $A+cuu^{\mathsf T}$ have equal rank, then
$u\in\operatorname{im}A$ and
$\operatorname{im}(A+cuu^{\mathsf T})=\operatorname{im}A$.
\end{lemma}
\begin{proof}
Put $d=\operatorname{rank}A$.  Symmetric elimination gives an
invertible matrix $S$ for which
\[
 S^{\mathsf T}AS=\begin{pmatrix}A_0&0\\0&0\end{pmatrix},
 \qquad \det A_0\ne0.
\]
Write $S^{\mathsf T}u=(v,w)^{\mathsf T}$.  If $u$ is outside
the image, some coordinate $w_j$ is nonzero.  The principal minor
on the $d$ coordinates of $A_0$ and this one further coordinate is
\begin{equation}\label{cr:minor}
 \det\begin{pmatrix}
 A_0+cvv^{\mathsf T}&cvw_j\\
 cw_jv^{\mathsf T}&cw_j^2
 \end{pmatrix}=cw_j^2\det A_0\ne0.
\end{equation}
For instance, subtract $(v_i/w_j)$ times the last row from row $i$
to verify the equality directly.  A rank-one update can increase
rank by at most one, proving the first claim.  Equal ranks therefore
force $u\in\operatorname{im}A$.  In that case the new image is
contained in the old one and has the same dimension, proving equality.
The proof remains valid when $w^{\mathsf T}w=0$; nonzero complex
isotropic vectors cause no exception to the argument.
\end{proof}

\begin{theorem}[Nonsingularity of rational symmetric primitives]
\label{cr:primitive}
Let $G(z)\in\operatorname{Sym}_q(F(z))$, let $w(z)\in F(z)$ be
nonzero, and let $p(z)\in F(z)^q$ have linearly independent
components over $F$.  If
\begin{equation}\label{cr:difference}
 G(z)-G(z-1)=w(z)p(z)p(z)^{\mathsf T},
\end{equation}
then $\det G(z)$ is a nonzero rational function.  This conclusion
holds for every symmetric rational solution, including after adding
an arbitrary constant symmetric matrix.
\end{theorem}
\begin{proof}
Suppose the generic rank were $d<q$.  A nonzero $d$-minor of $G$
has finitely many zeros and poles, and every $(d+1)$-minor vanishes
identically.  Hence there is an integer $N_0$ such that, at every
integer $N\ge N_0$, the matrices $G(N)$ and $G(N-1)$ are defined
and have rank $d$, while $w(N)\ne0$ and $p(N)$ is defined.
Characteristic zero ensures that these are infinitely many distinct
evaluation points.  Lemma~\ref{cr:rankone} applied to
\eqref{cr:difference} shows that the images of all these matrices
are one fixed $d$-dimensional subspace $W\subset F^q$, and that
$p(N)\in W$ at every sufficiently large integer.

Choose a nonzero linear functional $\lambda\in(F^q)^*$ vanishing
on $W$.  Then the rational function $\lambda p(z)$ vanishes at
infinitely many distinct points, so it is identically zero.  This
contradicts independence of the components of $p$.  Thus $d=q$.
Adding a constant symmetric matrix preserves \eqref{cr:difference},
so the same proof applies to every such addition.
\end{proof}

The result also applies with any fixed nonzero translation $h$ in
place of $1$, by changing variables.  Rationality is used to make
the ranks constant outside a finite set and to turn infinitely many
zeros into a function identity.  The statement is consequently stronger
than a positivity argument, but it does not make a claim about arbitrary
nonrational sequences of bilinear Gram matrices.

\subsection{The complex harmonic-degree theorem}

Let $C_1,\ldots,C_t\in\operatorname{Sym}_r(F)$ and put
\[
 C(X)=\sum_{j=1}^tX_jC_j,\qquad
 \rho=\operatorname{rank}_{F(X)}C(X),\qquad q=r-\rho.
\]
Suppose $R\in F(z)$ is nonzero and that the finite moment matrices
admit a decomposition
\begin{equation}\label{cr:decomposition}
 M_N=\sum_{n=0}^N R(n)v(n)v(n)^{\mathsf T}
     =D(N)+\sum_{j=1}^t h_j(N)C_j,
 \qquad v(n)=(1,n,\ldots,n^{r-1})^{\mathsf T},
\end{equation}
where $D(z)\in\operatorname{Sym}_r(F(z))$.  It suffices that the
decomposition and all its terms be defined for every sufficiently
large integer $N$.  The values $h_j(N)$ may lie in an extension of
$F$.  The variables $X_j$ are formal and independent; no arithmetic
independence of the actual harmonic values is asserted.

\begin{theorem}[Complex rational degree--rank equality]
\label{cr:degree}
Under these hypotheses,
\begin{equation}\label{cr:main}
 \boxed{\deg_X\det(D(z)+C(X))=\rho\quad\text{over }F(z).}
\end{equation}
Only finitely many admissible integer cutoffs can have smaller
degree.  In particular the conclusion holds for every nonzero complex
rational weight with no poles at nonnegative integers and its usual
finite harmonic decomposition.
\end{theorem}
\begin{proof}
The determinant has degree at most $\rho$, since any term choosing
more than $\rho$ columns from $C(X)$ is zero.  Since $F$ is infinite,
choose $a\in F^t$ for which $A=C(a)$ has rank $\rho$.  Let $K$
be an $r\times q$ constant matrix whose columns form its kernel.
Choose a complement $U$ so that $S=[U\ K]$ is invertible and
\[
 S^{\mathsf T}AS=\begin{pmatrix}A_0&0\\0&0\end{pmatrix},
 \qquad\det A_0\ne0.
\]
Every matrix $A+tC_j$ has rank at most $\rho$.  Over $F(t)$
its upper-left block is invertible near $t=0$, so its Schur complement
vanishes identically.  The coefficient of $t$ in that complement is
$K^{\mathsf T}C_jK$.  Therefore
\begin{equation}\label{cr:isotropy}
 K^{\mathsf T}C_jK=0\qquad(1\le j\le t).
\end{equation}
This algebraic argument is the complex form of the source's
compression lemma.

Put $G(z)=K^{\mathsf T}D(z)K$ and $p(z)=K^{\mathsf T}v(z)$.
The components of $p$ are independent polynomials, because $K$
has column rank $q$.  Equations \eqref{cr:decomposition} and
\eqref{cr:isotropy} show, at all sufficiently large integers, that
\[
 G(N)-G(N-1)=R(N)p(N)p(N)^{\mathsf T}.
\]
Both sides are rational functions in $N$, so this is the identity
\eqref{cr:difference}.  Theorem~\ref{cr:primitive} gives
$\det G(z)\ne0$.  Finally,
\begin{equation}\label{cr:direction}
 [t^\rho]\det(D(z)+tA)
   =\frac{\det A_0}{(\det S)^2}\det G(z)\ne0,
\end{equation}
obtained by selecting $tA_0$ in its $\rho$ rows and columns after
congruence by $S$.  This proves \eqref{cr:main}.  The same
coefficient is nonzero at all but finitely many admissible integer
cutoffs, proving the final assertion.  When $q=0$, its determinant
is one and the conclusion is immediate.
\end{proof}

No common kernel of the entire pencil is assumed.  The kernel of
one maximal-rank member is only totally isotropic for every $C_j$,
which is exactly the compression property needed in the proof.
The theorem answers the complex-weight question affirmatively rather
than restricting it to pencils with a common kernel.

\subsection{Finite certificates for every exceptional cutoff}

The proof gives a cheap sufficient certificate using just one
maximal-rank direction.  Choose a nonzero polynomial $d(z)$ clearing
the denominators of $G(z)$ and set $H(z)=d(z)G(z)$.  Then
\begin{equation}\label{cr:certificate}
 0\ne E(z):=\det H(z)\in F[z].
\end{equation}
Every degree-deficient cutoff at which $d(N)\ne0$ is a root of
$E$.  If $e_{ij}=\deg H_{ij}$, with degree $-\infty$ for zero
entries, then
\begin{equation}\label{cr:degree-bound}
 \deg E\le \max_{\pi\in S_q}\sum_{i=1}^q e_{i,\pi(i)}
 \le q\max_{i,j}e_{ij}.
\end{equation}
Thus the right side bounds the number of distinct such exceptions.
This is a bound computable from rational entries and does not depend
on ordering the coefficient field.  Apparent poles introduced by
denominator clearing must be removed or checked in the original
decomposition.

For an exact rather than sufficient certificate, write the highest
homogeneous component after clearing a common denominator in $z$ as
\[
 [\det(D(z)+C(X))]_\rho
       =\frac{1}{b(z)}\sum_{|\nu|=\rho}e_\nu(z)X^\nu.
\]
Theorem~\ref{cr:degree} guarantees that not all $e_\nu$ vanish.
For an admissible $N$ with $b(N)\ne0$, degree is deficient exactly
when $g(N)=0$, where
\begin{equation}\label{cr:gcd}
 g(z)=\gcd_{|\nu|=\rho}e_\nu(z).
\end{equation}
This is the complex-field extension
of the source's coefficient-gcd test; it is not a new polynomial-gcd
algorithm.  The cheaper polynomial $E$ supplies a finite candidate
set, and the full highest component filters it exactly.

\subsection{A residue interpretation of the rational compression}

For clarity, the usual harmonic decomposition is available over
$\mathbb C$ without a sign assumption.  Partial fractions of
$z^mR(z)$ consist of a polynomial and finitely many terms
$c/(z-\alpha)^j$.  Poles differing by integers belong to one shift
orbit.  Summing those terms from $0$ to $N$ produces one generalized
harmonic coordinate per orbit and pole order, together with rational
endpoint corrections.  This is the form \eqref{cr:decomposition}.

The discrete residue of a rational function at an orbit and pole
order is the sum of the corresponding partial-fraction coefficients
over that orbit.  The classical rational-summability criterion states
that a rational function has a rational antidifference if and only if
all these discrete residues vanish \cite[Proposition~2.5]{ChenSinger2012}.
For a polynomial subspace $V$ with basis $p_1,\ldots,p_q$, the
following statements are equivalent:
\begin{enumerate}
\item every $R(z)p_i(z)p_j(z)$ has a rational antidifference;
\item every discrete-residue bilinear form
$(f,g)\mapsto\operatorname{dres}_{\mathcal O,k}(Rfg)$ vanishes
on $V\times V$;
\item if $K$ is the coefficient matrix of this basis, then
$K^{\mathsf T}C_jK=0$ for every harmonic coefficient matrix.
\end{enumerate}
The equivalence is the classical residue criterion applied entry by
entry.  Theorem~\ref{cr:primitive} adds the following consequence:
whenever these conditions hold and the $p_i$ are independent, every
symmetric rational matrix primitive of the products $Rp_ip_j$ is
nonsingular over $\mathbb C(z)$.  Thus a maximal-rank pencil kernel
produces a rationally summable polynomial space whose accumulated
bilinear form cannot be identically degenerate.

\subsection{Exact nonreal examples}

Let
\[
 R_\lambda(z)=\frac1{z+1}+\frac\lambda{z+2},\qquad\lambda\ne0,
 \qquad X=\sum_{n=0}^N\frac1{n+1}.
\]
For an $r\times r$ moment matrix, $r\ge2$, the only harmonic
coefficient matrix is
\[
 C=u_1u_1^{\mathsf T}+\lambda u_2u_2^{\mathsf T},\qquad
 u_a=(1,-a,\ldots,(-a)^{r-1})^{\mathsf T},
\]
and has rank two.  Thus the harmonic degree is exactly two over
$\mathbb C(N)$, for every nonzero complex $\lambda$.

At $r=3$, the kernel polynomial is $(z+1)(z+2)$ and
\begin{equation}\label{cr:three}
 \boxed{
 [X^2]\det(D(N)+XC)=
 \frac{\lambda}{12}(N+1)(N+2)(N+3)
       \{3(\lambda+1)N+4\lambda+8\}.}
\end{equation}
At $r=4$, use the two kernel polynomials
$(z+1)(z+2)$ and $z(z+1)(z+2)$.  Then
\begin{equation}\label{cr:four}
 \begin{split}
 [X^2]\det(D(N)+XC)
 ={}&\frac{\lambda}{3600}N(N+1)^2(N+2)^2(N+3)^2(N+4)\\
 &\quad\cdot\bigl\{6(\lambda+1)^2N^2
 +(14\lambda^2+48\lambda+34)N
 +10\lambda^2+45\lambda+50\bigr\}.
 \end{split}
\end{equation}
Both identities follow from \eqref{cr:direction} and finite sums of
polynomials.  In the four-dimensional case, the two-by-two compressed
determinant is the expression after the initial factor $\lambda$.
For $\lambda=i$, its final quadratic is
$20N+40+i(12N^2+48N+45)$, so it cannot vanish at a nonnegative
integer.  Consequently the four-dimensional harmonic degree is
exactly two at every $N\ge1$, even though the actual evaluated
moment determinant vanishes at $N=1,2$ because too few nodes are
present.  Formal degree and numerical determinant value remain
different notions.

A nonreal rational weight can nevertheless have a genuine later
degree exception.  Take $r=4$ and
\begin{equation}\label{cr:exception-weight}
 R(z)=\frac{39}{z+1}+\frac{78i}{z+2}
                        -\frac{61+95i}{z+3}.
\end{equation}
Its one harmonic matrix has rank three.  The kernel polynomial is
$p(z)=(z+1)(z+2)(z+3)$, and exact Faulhaber summation gives
\begin{equation}\label{cr:exception-compression}
 \begin{split}
 \sum_{n=0}^N R(n)p(n)^2
  =-\frac16&(N-3)(N+1)(N+2)(N+3)(N+4)\\
           &\cdot\{(22+17i)N+56+22i\}.
 \end{split}
\end{equation}
The top-left three-square minor of the harmonic matrix is
$1155960-742248i\ne0$.  Hence the only degree-deficient
nonnegative cutoff is $N=3$.  At that cutoff the formal determinant
is exactly
\begin{equation}\label{cr:exception-polynomial}
 \begin{split}
 P_3(X)={}&(-22046418576+2049048144i)X^2\\
 &+\left(\frac{455872142232}{5}-7792299840i\right)X\\
 &-\frac{2356478577592}{25}
                         +\frac{183512850456}{25}i.
 \end{split}
\end{equation}
Thus its degree drops from three to two, while substitution of the
actual $X=25/12$ gives
\[
 P_3(25/12)=\frac{55033-146419i}{25}\ne0.
\]
Independently, this value is the four-node Vandermonde square times
$R(0)R(1)R(2)R(3)$.  The exception is therefore a formal-degree
phenomenon at a nonzero evaluated determinant, and not a lack of
nodes or a vanishing weight.

\subsection{Integration and further questions}

The source's real-weight results are valid and need no retraction.
Its open complex-weight question can be marked solved under the
same rational-decomposition hypotheses, and Corollary
\texttt{deg:rational} can be broadened to complex rational weights.
The eventual-sign bound $r(r-\rho)$ should retain its original real
hypotheses.  The present denominator-degree certificate is a
different bound and is not asserted to equal that real bound.

Two precise next questions remain.  First, determine sharp bounds
for the exceptional-cutoff polynomial using only the pole orders and
integer dispersions of $R$, without constructing the whole compressed
primitive.  Second, classify polynomial subspaces totally isotropic
for all discrete-residue forms of $Rfg$; the primitive theorem proves
nonsingularity once such a subspace is given, but does not classify
the subspaces or the rank-two pencils producing them.


% END sections/02_complex_rank.tex

% BEGIN sections/03_hahn_and_polylogs.tex
\section{Adjacent poles, Hahn determinants, and classical polylogarithms}
\label{hp:section}

The complex degree theorem explains why a rational weight may have
surprisingly small harmonic complexity even in a large determinant.
We now make that statement explicit for two adjacent simple poles.
The general Christoffel determinant transformation used below is
classical; its particularly compact form and its harmonic consequences
are the specialization developed here. We use the formulation in
\cite{Krattenthaler2023} and give the short determinant proof needed
for our normalization. The finite orthogonal polynomials are the
uniform Hahn, or discrete Chebyshev, polynomials
\cite[Sections 18.19--18.20]{DLMF}.

\subsection{A three-by-three formula for every particle count}

For $r\ge2$ and an admissible shift $a$, set
\begin{align}
 R_\eta(n;a)&=\frac1{n+a}+\frac{\eta}{n+a+1},\notag\\
 Z_{r,N}(a,\eta)
 &=\det_{0\le i,j<r}
       \left(\sum_{n=0}^N n^{i+j}R_\eta(n;a)\right).
 \label{hp:ensemble}
\end{align}
Admissible means that no displayed denominator vanishes. For
$N\ge r-1$, finite Cauchy--Binet gives the equivalent ensemble
\begin{equation}\label{hp:vandermonde}
 Z_{r,N}(a,\eta)=
 \sum_{0\le n_1<\cdots<n_r\le N}
       \prod_{i<j}(n_j-n_i)^2\prod_{\ell=1}^rR_\eta(n_\ell;a).
\end{equation}
An empty sum is zero. Resolve the shifted harmonic coordinate by
\[
 X=H_N^{(1)}(a):=\sum_{n=0}^N\frac1{n+a},
 \qquad
 \sum_{n=0}^N\frac1{n+a+1}
 =X+\frac1{N+a+1}-\frac1a.
\]
The determinant becomes a formal polynomial
$\mathcal P_{r,N}(X;a,\eta)$ with rational coefficients in the
cutoff and parameters. For $\eta\ne0$ its harmonic pencil has rank
two, so its generic degree in $X$ is exactly two. Write
$T_{r,N}(a,\eta)=[X^2]\mathcal P_{r,N}$.

\begin{proposition}[A polynomially weighted compression]
\label{hp:compression}
Put $q=r-2$ and $L_a(n)=(n+a)(n+a+1)$. Then
\begin{equation}\label{hp:top-compressed}
 \boxed{\displaystyle
 T_{r,N}(a,\eta)=\eta
 \det_{0\le i,j<q}\left(
 \sum_{n=0}^N n^{i+j}L_a(n)
                   \{(1+\eta)(n+a)+1\}\right).}
\end{equation}
The determinant of size zero is one.
\end{proposition}
\begin{proof}
The harmonic pencil is $C=uu^{\mathsf T}+\eta vv^{\mathsf T}$,
where $u_i=(-a)^i$ and $v_i=(-a-1)^i$. Its common kernel consists
of coefficient vectors of polynomials divisible by $L_a$. Use the
basis
\[
 1,\ n,\ L_a(n),\ nL_a(n),\ldots,n^{q-1}L_a(n).
\]
Its monic triangular change of basis has determinant one. The
upper-left two-by-two minor of $C$ has determinant $\eta$, while
the lower-right compression is a sum of polynomial weights:
\[
 R_\eta(n;a)L_a(n)^2
   =L_a(n)\{(1+\eta)(n+a)+1\}.
\]
Selecting the two harmonic entries in the transformed determinant
proves \eqref{hp:top-compressed}.
\end{proof}

Define the monic uniform Hahn polynomial by the finite expression
\begin{equation}\label{hp:hahn}
 p_k(x;N)=\binom{2k}{k}^{-1}
 \sum_{j=0}^k(-1)^{k+j}\binom kj\binom{k+j}{j}
                  (N-j)^{\underline{k-j}}x^{\underline j}.
\end{equation}
For $0\le k\le N$ it is orthogonal on $\{0,\ldots,N\}$ with
unit weights. Its squared norm is
\begin{equation}\label{hp:norm}
 h_k(N)=\sum_{n=0}^Np_k(n;N)^2
   =\frac{(N+k+1)!}{(2k+1)(N-k)!\binom{2k}{k}^2},
 \qquad \mathcal H_q(N)=\prod_{k=0}^{q-1}h_k(N).
\end{equation}
The finite formula also avoids apparent singularities in the
usual terminating hypergeometric notation
\[
 p_k(x;N)=
 \frac{(-1)^kN^{\underline k}}{\binom{2k}{k}}
 {}_3F_2\!\left(\begin{matrix}-k,k+1,-x\\1,-N\end{matrix};1\right).
\]

\begin{theorem}[A fixed-size Hahn formula]\label{hp:hahn-top}
Let $r\ge2$, $N\ge r$, $\eta\ne-1$, and put $t=1+\eta$ and
\[
 c_1=-a,\qquad c_2=-a-1,\qquad c_3=-a-\frac1t,\qquad q=r-2.
\]
Then
\begin{equation}\label{hp:hahn-top-formula}
 \boxed{\displaystyle
 T_{r,N}(a,\eta)
   =(-1)^q t^r\mathcal H_q(N)
           \det_{1\le i\le3,\ 0\le j\le2}
                         p_{q+j}(c_i;N).}
\end{equation}
The expression has a removable limit at $\eta=-1$ and a zero at
$\eta=0$. The compressed formula \eqref{hp:top-compressed}
specifies those cases without a limit. Thus a determinant of
arbitrary size has a highest harmonic coefficient represented by
one determinant of size three.
\end{theorem}
\begin{proof}
For monic orthogonal polynomials $p_k$ of a measure $\nu$, let
$D_q[\nu]$ be its moment determinant. For distinct $c_1,\ldots,c_d$,
the Christoffel transformation is
\begin{equation}\label{hp:christoffel}
 D_q\!\left[\prod_{\ell=1}^d(x-c_\ell)\nu\right]
 =(-1)^{dq}D_q[\nu]\,
       \frac{\det(p_{q+j-1}(c_i))_{i,j=1}^d}
            {\prod_{i<j}(c_j-c_i)}.
\end{equation}
One way to verify the sign and normalization is to append the
$d$ fixed nodes $c_i$ to the $q$ integration nodes in the
Vandermonde determinant, replace its monomial columns by the
monic orthogonal columns, and integrate the first $q$ rows
against their companion Vandermonde. Orthogonality leaves the
product of the first $q$ squared norms; the remaining $d$ rows
give the displayed determinant. The cross factors are
$\prod_{\ell,j}(c_\ell-x_j)$, producing $(-1)^{dq}$.
This also proves \eqref{hp:christoffel} for a finite discrete
measure by finite Cauchy--Binet.

In \eqref{hp:top-compressed}, the weight is
$t\prod_{\ell=1}^3(n-c_\ell)$. Moreover
\[
 \prod_{i<j}(c_j-c_i)=\frac{\eta}{t^2}.
\]
Substitution in \eqref{hp:christoffel}, followed by multiplication
by the exterior $\eta$, gives
$(-1)^{3q}t^{q+2}\mathcal H_q\det p$, which is
\eqref{hp:hahn-top-formula}. The proof initially assumes
$\eta\ne0,-1$; the extension follows from the polynomial
identity \eqref{hp:top-compressed}.
\end{proof}

\subsection{Every exceptional residue parameter is an orthogonal zero}

The complex degree theorem allows a finite exceptional set of
cutoffs for a fixed weight. Here one can solve the dual question:
for a fixed cutoff, which pole residues cause degree loss?

\begin{corollary}[Exact exceptional parameters]\label{hp:roots}
Let $a>0$, $q=r-2\ge1$, and $N\ge q$. Let $\pi_q$ be the monic
orthogonal polynomial of degree $q$ for the positive discrete
measure with weights $L_a(n)$ on $\{0,\ldots,N\}$, and let
$\xi_1,\ldots,\xi_q$ be its zeros. These zeros are distinct and
belong to $(0,N)$. Put
\[
 \mathcal H_q[L_a]=
 \det_{0\le i,j<q}\left(\sum_{n=0}^N n^{i+j}L_a(n)\right)>0.
\]
Then
\begin{equation}\label{hp:root-factorization}
 T_{r,N}(a,\eta)
  =\eta\,\mathcal H_q[L_a]\,
         \prod_{j=1}^q\{(1+\eta)(a+\xi_j)+1\}.
\end{equation}
Apart from $\eta=0$, where the pencil rank changes, all
degree-deficient parameters are exactly
\begin{equation}\label{hp:exceptional}
             \eta_j=-1-\frac1{a+\xi_j}\quad(1\le j\le q).
\end{equation}
They are simple, real, and less than $-1$. In particular no
nonreal residue parameter causes degree loss at these cutoffs.
\end{corollary}
\begin{proof}
The one-factor case of \eqref{hp:christoffel}, with
$c=-a-1/(1+\eta)$, gives
$D_q[((1+\eta)(n+a)+1)L_a]
=\mathcal H_q[L_a]\prod_j\{(1+\eta)(a+\xi_j)+1\}$.
The standard zero theorem for a positive orthogonal measure
gives simplicity and the location \cite[Section 18.2(vi)]{DLMF}.
Equation \eqref{hp:top-compressed} finishes the proof.
\end{proof}

These statements concern formal harmonic degree. They do not
assert that the evaluated determinant has a zero, and they do
not specify how far its degree drops beyond the vanishing of
the quadratic term. The nonreal example in
Section~\ref{cr:section} shows why this distinction matters.

\subsection{A uniform classical-polylogarithm formula}

Now specialize to $a=1$ and use $M=N+1$. Passing from monomials
in $n$ to monomials in $k=n+1$ preserves the determinant. Write
$S_j(M)=\sum_{k=1}^M k^j$, regarded as its Faulhaber polynomial,
and let $X$ be the formal replacement for $H_M$. The moment of
degree $j\ge0$ is
\begin{equation}\label{hp:compiler}
 \mu_j(M,X;\eta)=
 \delta_{j0}X+\mathbf1_{j\ge1}S_{j-1}(M)
 +\eta\left\{
   \sum_{\ell=0}^{j-1}(-1)^{j-1-\ell}S_\ell(M)
   +(-1)^j\left(X+\frac1{M+1}-1\right)\right\}.
\end{equation}
The sum is empty at $j=0$. This is simply polynomial division
of $k^j/(k+1)$. Consequently
\[
 \mathcal P_r(M,X;\eta)=\det(\mu_{i+j})_{0\le i,j<r}
\]
is an explicit finite construction of every cutoff formula.

\begin{theorem}[Only a dilogarithm is needed for every $r$]
\label{hp:all-r-generator}
For $r\ge2$ there are polynomials $A_{rj}(M;\eta)$,
$0\le j\le2$, such that
\begin{equation}\label{hp:polynomial-plus-residue}
 \mathcal P_r(M,X;\eta)=\sum_{j=0}^2A_{rj}(M;\eta)X^j
                        +\frac{B_r(X;\eta)}{M+1},
\end{equation}
where the entire endpoint residue is
\begin{equation}\label{hp:residue}
 B_2=\eta(X-\eta-2),\qquad B_3=-\eta,\qquad B_r=0\ (r\ge4).
\end{equation}
Let $\theta=x\,d/dx$, $L=-\log(1-x)$, and
\[
 F_0(x)=\frac{x}{1-x},\qquad F_1(x)=\frac{L}{1-x},
 \qquad F_2(x)=\frac{\operatorname{Li}_2(x)+L^2}{1-x}.
\]
For $|x|<1$ the cutoff generating function is
\begin{equation}\label{hp:generator-general}
 \boxed{\displaystyle
 \sum_{M\ge1}Z_{r,M-1}(1,\eta)x^M
       =\sum_{j=0}^2A_{rj}(\theta;\eta)F_j(x)+\mathcal R_r(x;\eta),}
\end{equation}
with
\begin{align}
 \mathcal R_2&=\eta\left\{\frac{L^2}{2x}
                   -(\eta+2)\left(\frac Lx-1\right)\right\},\notag\\
 \mathcal R_3&=-\eta\left(\frac Lx-1\right),\qquad
 \mathcal R_r=0\quad(r\ge4).
 \label{hp:generator-residue}
\end{align}
All apparent singularities at $x=0$ are removable.
For every particle count the result therefore belongs to the
space over $\mathbb C(x)$ generated by
$1,L,L^2,\operatorname{Li}_2(x)$.
\end{theorem}
\begin{proof}
In the $k$ basis put $e_0=(1,0,\ldots)^{\mathsf T}$,
$e_1=(0,1,0,\ldots)^{\mathsf T}$, and
$v_i=(-1)^i$. The matrix in \eqref{hp:compiler} has the form
\[
 A(M)+X(e_0e_0^{\mathsf T}+\eta vv^{\mathsf T})
                       +\frac{\eta vv^{\mathsf T}}{M+1},
\]
where $A$ is polynomial in $M$. The harmonic part has rank at
most two, and the pole has rank one. The determinant thus has
degree at most two in $X$ and at most a simple pole at $M=-1$.

The polynomial identities $S_0(-1)=-1$ and $S_j(-1)=0$ for
$j\ge1$ show that the regular matrix at that point is
\[
 X(e_0e_0^{\mathsf T}+\eta vv^{\mathsf T})
              -\eta e_0e_0^{\mathsf T}
              -(e_0e_1^{\mathsf T}+e_1e_0^{\mathsf T}).
\]
Its residue after the rank-one update is directly
$\eta(X-\eta-2)$ at $r=2$ and $-\eta$ at $r=3$.
For $r\ge4$, the matrix even after adding an arbitrary multiple
of $vv^{\mathsf T}$ has image contained in
$\operatorname{span}(e_0,e_1,v)$, of dimension at most three.
Its determinant is identically zero, so the pole residue
vanishes. Polynomial division proves
\eqref{hp:polynomial-plus-residue}.

The elementary harmonic generating identities are
\[
 \sum_{M\ge1}H_Mx^M=F_1(x),\qquad
 \sum_{M\ge1}H_M^2x^M=F_2(x).
\]
For the second, split $H_M^2$ into its diagonal terms and twice
its strict off-diagonal terms, then sum the outer cutoff.
Multiplication by a polynomial in $M$ becomes its substitution
at $\theta$. Finally
\[
 \sum_{M\ge1}\frac{x^M}{M+1}=\frac Lx-1,\qquad
 \sum_{M\ge1}\frac{H_Mx^M}{M+1}=\frac{L^2}{2x};
\]
the latter follows by integrating $F_1$. These facts prove the
formula and convergence. Differentiation by $\theta$ preserves
the displayed four-function space, since
$d\operatorname{Li}_2(x)/dx=L/x$.
\end{proof}

This is a closed finite recipe, with polynomial coefficient
construction \eqref{hp:compiler}. It does not claim that the
four functions are always a minimal basis. Special parameters
can produce further cancellations. Nor does formal harmonic
degree in a general problem automatically equal minimal
multiple-polylogarithm depth; the conclusion here uses the
specific pole cancellation proved above.

\subsection{A positive three-particle harmonic identity}

The choice $\eta=-1$ makes the weight positive and especially
simple:
\[
 R_{-1}(n;1)=\frac1{(n+1)(n+2)}.
\]
Let
\[
 \mathcal Z_N=
 \sum_{0\le n_1<n_2<n_3\le N}
   \frac{\prod_{i<j}(n_j-n_i)^2}
        {\prod_{i=1}^3(n_i+1)(n_i+2)}.
\]
The first two values vanish. Direct evaluation of the three-by-three
moment determinant in \eqref{hp:compiler} gives the following identity
for every integer $M\ge1$:
\begin{equation}\label{hp:triple}
 \boxed{\displaystyle
 \mathcal Z_{M-1}
 =-\frac{M(M+1)(M+2)}3H_M^2
   +\frac{4M^2(M+2)}3H_M
   +\frac{M^3(M^2-18M-31)}{12(M+1)}.}
\end{equation}
For example $\mathcal Z_2=1/36$ and $\mathcal Z_3=23/90$.
Equation \eqref{hp:triple} is also an exact digamma formula
after $H_M=\psi(M+1)+\gamma$.

Its entire cutoff generator has a short classical expression.
With $B=\operatorname{Li}_2(x)+L^2$,
\begin{equation}\label{hp:triple-generator}
 \boxed{\begin{aligned}
 G(x):=\sum_{M\ge1}\mathcal Z_{M-1}x^M
 ={}&-\frac{2xB}{(1-x)^4}
  +\frac{(3-14x+15x^2-2x^3)L}{3x(1-x)^4}\\
 &+\frac{18x^3-31x^2+31x-6}{6(1-x)^5}.
 \end{aligned}}
\end{equation}
The right side starts $x^3/36+23x^4/90+\cdots$ despite its
apparent low-order terms. One proof is to apply
Theorem~\ref{hp:all-r-generator} to \eqref{hp:triple}; it reduces
the statement to differentiating $F_0,F_1,F_2$ a finite number
of times. Thus both the sum and the displayed closed form
are established without a numerical relation search.

\subsection{A convergent zeta-valued endpoint}

The determinant is $O(M^4+M^3\log^2 M)$, so its sixth Euler
primitive has an ordinary finite value at one.

\begin{corollary}[An explicit ensemble integral]\label{hp:endpoint}
\begin{align}
 \mathcal E
 &:=\sum_{M\ge1}\frac{\mathcal Z_{M-1}}{M^6}
   =\frac1{120}\int_0^1\frac{(-\log x)^5}{x}G(x)\,dx
       \notag\\
 &=-1+\frac{13}{12}\zeta(2)-\frac{31}{12}\zeta(3)
       +\frac53\zeta(4)+\frac{41}{6}\zeta(5)
       -\frac{97}{24}\zeta(6)-4\zeta(7)\notag\\
 &\hspace{1em}-\frac73\zeta(2)\zeta(3)+2\zeta(3)^2
       +\frac23\zeta(2)\zeta(5)+\frac53\zeta(3)\zeta(4).
 \label{hp:endpoint-value}
\end{align}
Numerically $\mathcal E=0.0027314671316270802642692390978\ldots$.
\end{corollary}
\begin{proof}
The elementary integral
$\int_0^1x^{M-1}(-\log x)^5\,dx=5!/M^6$ and positivity
justify interchange by Tonelli. Set
$T_s=\sum_{M\ge1}H_M/M^s$ and
$S_s=\sum_{M\ge1}H_M^2/M^s$. Equation \eqref{hp:triple}
and the partial fraction
\[
 \frac{M^2-18M-31}{12M^3(M+1)}
 =\frac1{M+1}-\frac1M+\frac{13}{12M^2}-\frac{31}{12M^3}
\]
give
\begin{equation}\label{hp:euler-reduction}
 \mathcal E=-\frac{S_3+3S_4+2S_5}{3}
         +\frac43(T_3+2T_4)
         -1+\frac{13}{12}\zeta(2)-\frac{31}{12}\zeta(3).
\end{equation}
The classical Euler-sum reductions required here are
\begin{align}
 T_3&=\tfrac54\zeta(4),&
 T_4&=3\zeta(5)-\zeta(2)\zeta(3),\notag\\
 S_3&=\tfrac72\zeta(5)-\zeta(2)\zeta(3),&
 S_4&=\tfrac{97}{24}\zeta(6)-2\zeta(3)^2,\notag\\
 S_5&=6\zeta(7)-\zeta(2)\zeta(5)
                         -\tfrac52\zeta(3)\zeta(4).
 \label{hp:euler-values}
\end{align}
Substitution proves the result. For completeness, the next
paragraph derives these classical inputs, so the endpoint
evaluation does not depend on a table of numerical relations.
\end{proof}

Write
$\zeta(s_1,\ldots,s_d)=\sum_{k_1>\cdots>k_d\ge1}
k_1^{-s_1}\cdots k_d^{-s_d}$. Directly sorting the harmonic
indices gives
\begin{align}
 T_s&=\zeta(s,1)+\zeta(s+1),\notag\\
 S_s&=2\zeta(s,1,1)+\zeta(s,2)
                       +2\zeta(s+1,1)+\zeta(s+2).
 \label{hp:sorting}
\end{align}
These are particular cases of the established Euler-sum/MZV
calculus \cite{Xu2017,XuWang2019}. The height-one generating
identity, also classical, is
\begin{equation}\label{hp:height-one}
\begin{aligned}
 \sum_{n\ge2,\ m\ge0}\zeta(n,\{1\}^m)X^{n-1}Y^{m+1}
 &=1-\frac{\Gamma(1-X)\Gamma(1-Y)}{\Gamma(1-X-Y)}\\
 &=1-\exp\!\sum_{k\ge2}\frac{\zeta(k)}k
                   \{X^k+Y^k-(X+Y)^k\}.
\end{aligned}
\end{equation}
A self-contained proof starts from
$\sum_mH_{N-1}(\{1\}^m)Y^m=\prod_{j<N}(1+Y/j)$.
Sum first over $n$ and then $N$, to obtain
\[
 Y\int_0^1(t^{-X}-1)(1-t)^{-1-Y}\,dt
 =Y\{B(1-X,-Y)-B(1,-Y)\}.
\]
Initially take $\Re Y<0$ and $X,Y$ small; the beta integral
gives the gamma quotient and analytic continuation gives the
identity at zero. Its coefficients yield Euler's formula
\[
 \zeta(n,1)=\frac n2\zeta(n+1)
        -\frac12\sum_{j=1}^{n-2}\zeta(j+1)\zeta(n-j)
\]
and, at the orders needed in \eqref{hp:sorting},
\begin{align*}
 \zeta(3,1,1)&=2\zeta(5)-\zeta(2)\zeta(3),\\
 \zeta(4,1,1)&=\tfrac{23}{16}\zeta(6)-\zeta(3)^2,\\
 \zeta(5,1,1)&=5\zeta(7)-2\zeta(2)\zeta(5)
                                     -\tfrac54\zeta(3)\zeta(4).
\end{align*}
The remaining double zeta values follow from the convergent
stuffle identity
$\zeta(a)\zeta(b)=\zeta(a,b)+\zeta(b,a)+\zeta(a+b)$
and the iterated-integral shuffle identity
\[
 \zeta(a)\zeta(b)=
 \sum_{j=0}^{a-1}\binom{b-1+j}{j}\zeta(b+j,a-j)
 +\sum_{j=0}^{b-1}\binom{a-1+j}{j}\zeta(a+j,b-j).
\]
Use $(a,b)=(2,3),(3,3),(2,5),(3,4)$ and the height-one
values just obtained. The resulting finite linear systems give
\begin{align*}
 \zeta(3,2)&=3\zeta(2)\zeta(3)-\tfrac{11}{2}\zeta(5),\\
 \zeta(4,2)&=\zeta(3)^2-\tfrac43\zeta(6),\\
 \zeta(5,2)&=5\zeta(2)\zeta(5)+2\zeta(3)\zeta(4)-11\zeta(7).
\end{align*}
Together with the elementary even-zeta products these prove
\eqref{hp:euler-values}. We claim the ensemble specialization
and its resulting integral as the additional identities here;
the Euler-sum values and the gamma generating identity themselves
are established background.


% END sections/03_hahn_and_polylogs.tex

% BEGIN sections/04_collisions.tex
% Integration fragment; all labels use the urc: prefix.
% Requires amsmath, amssymb, amsthm, hyperref. No custom global macros.
% Bibliography: cite the existing DLMF and Lattice Contacts entries.

\section{Unequal-rate Stieltjes boundary collisions}
\label{urc:section}

The independent-shift product generator in the incoming \emph{Lattice
Contacts and Mixed Zeta Identities} treats positive shifts and proves
compatibility with the coincident finite part along the common-shift
diagonal. Its Question~Q9, repeated at the end of its independent-shift
section \cite[Section~6 and Question~Q9]{LatticeIncoming}, leaves open the path
\[
 a=\lambda\epsilon,\qquad b=\mu\epsilon,\qquad
 \lambda,\mu>0,\qquad \epsilon\downarrow0.
\]
We determine every divergent power and logarithm on this path and its
remaining constant, for arbitrary Stieltjes indices and argument
derivatives. The constant correction is considerably smaller than the
full independent-shift product generator: it depends only on two Taylor
coefficients at one and a finite polynomial in the shift ratio.

These statements resolve the specified collision question relative to
the inspected repository sources. The underlying local-subtraction
method is classical; no claim of historical priority for general
finite-part asymptotics is intended.

\subsection{The coordinate and the finite contact polynomial}

We use the convention
\begin{equation}\label{urc:stieltjes}
 Z(u,z)=\zeta(1+u,z)-\frac1u
       =\sum_{m\ge0}\frac{(-1)^m}{m!}\gamma_m(z)u^m.
\end{equation}
In particular, $\gamma_0(z)=-\psi(z)$. Write
\begin{equation}\label{urc:integral}
 I_{m,n}^{r,t}(\epsilon;\lambda,\mu)
 =\int_0^1\gamma_m^{(r)}(x+\lambda\epsilon)
             \gamma_n^{(t)}(x+\mu\epsilon)\,dx,
 \qquad m,n,r,t\in\mathbb Z_{\ge0}.
\end{equation}
Every integral in \eqref{urc:integral} is ordinary. Its asymptotic
constant means the coefficient of $\epsilon^0(\log\epsilon)^0$.
The inherited coincident finite part is
\begin{equation}\label{urc:finitepart}
 F_{m,n}^{r,t}
 =\operatorname{CT}_{\delta\downarrow0}
    \int_\delta^1\gamma_m^{(r)}(x)\gamma_n^{(t)}(x)\,dx.
\end{equation}
Thus the two constants use the displayed coordinates $\epsilon$ and
$\delta$; no subtraction scale is left implicit.

For $r\ge0$ define a holomorphic germ at $u=0$ by
\begin{equation}\label{urc:contact-germ}
 C_r(u;\lambda,\mu)=\frac{(-1)^r(1+u)_r}{r!}
 \left\{\frac{\lambda^{-u}-1}{u}
 +\lambda^{-u}\sum_{j=1}^r
    \binom rj\frac{(\mu/\lambda-1)^j}{u+j}\right\},
\end{equation}
and put
\begin{equation}\label{urc:kappa}
 \kappa_{m,r}(\lambda,\mu)
 =(-1)^m\left.\partial_u^m C_r(u;\lambda,\mu)\right|_{u=0}.
\end{equation}
The sum is empty if $r=0$. The quotient at $u=0$ is removable.
Moreover, $(1+u)_r/(u+j)$ is a polynomial for $1\le j\le r$,
so \eqref{urc:contact-germ} contains no hidden special-function limit.

\begin{theorem}[Complete constant correction for unequal rates]
\label{urc:constant-theorem}
For all indices in \eqref{urc:integral},
\begin{equation}\label{urc:constant-formula}
 \boxed{\begin{aligned}
 \operatorname{CT}_{\epsilon\downarrow0}
       I_{m,n}^{r,t}(\epsilon;\lambda,\mu)
 ={}&F_{m,n}^{r,t}
    +\gamma_n^{(r+t)}(1)\kappa_{m,r}(\lambda,\mu)\\
   &+\gamma_m^{(r+t)}(1)\kappa_{n,t}(\mu,\lambda).
 \end{aligned}}
\end{equation}
In particular, with $R=\mu/\lambda$,
\begin{equation}\label{urc:zeroth-contact}
 \boxed{\displaystyle
 \kappa_{0,r}(\lambda,\mu)
 =(-1)^r\left\{-\log\lambda+
                   \sum_{j=1}^r\frac{R^j-1}{j}\right\}.}
\end{equation}
For arbitrary Stieltjes indices without argument differentiation,
\begin{equation}\label{urc:no-argument-derivatives}
 \operatorname{CT}_{\epsilon\downarrow0}I_{m,n}^{0,0}
 =F_{m,n}^{0,0}
  -\frac{\gamma_n(1)}{m+1}(\log\lambda)^{m+1}
  -\frac{\gamma_m(1)}{n+1}(\log\mu)^{n+1}.
\end{equation}
\end{theorem}

Only the Taylor coefficient of degree $r$ in the analytic part of
the second factor and the coefficient of degree $t$ in the analytic
part of the first factor enter \eqref{urc:constant-formula}. In
particular, the singular--singular product changes the divergent
terms but contributes no path-dependent constant.

\subsection{A finite formula for every contact coefficient}

For a fully finite harmonic description, write
\[
 A_r(u)=\frac{(1+u)_r}{r!}
       =\sum_{a=0}^r e_{r,a}u^a,
 \qquad
 Q_r(u;R)=\sum_{j=1}^r\binom rj(R-1)^j\frac{A_r(u)}{u+j}
       =\sum_{a=0}^{r-1}q_{r,a}(R)u^a.
\]
The $Q_0$ sum is empty. The coefficients $e_{r,a}$ are elementary
symmetric polynomials in $1,1/2,\ldots,1/r$, equivalently
\begin{equation}\label{urc:harmonic-generator}
 A_r(u)=\exp\left\{\sum_{h\ge1}
        \frac{(-1)^{h+1}}h H_r^{(h)}u^h\right\}.
\end{equation}
With $L=\log\lambda$, coefficient extraction in
\eqref{urc:contact-germ} gives
\begin{equation}\label{urc:finite-coefficients}
 \begin{aligned}
 \kappa_{m,r}={}&(-1)^{m+r}m!
 \sum_{a=0}^{\min(r,m)}e_{r,a}
          \frac{(-L)^{m-a+1}}{(m-a+1)!}\\
 &+(-1)^{m+r}m!
 \sum_{a=0}^{\min(r-1,m)}q_{r,a}(R)
          \frac{(-L)^{m-a}}{(m-a)!}.
 \end{aligned}
\end{equation}
The second sum is interpreted as empty at $r=0$. For example,
\begin{align*}
 \kappa_{0,0}&=-L,&
 \kappa_{1,0}&=-\frac{L^2}{2},\\
 \kappa_{0,1}&=L+1-R,&
 \kappa_{1,1}&=\frac{L^2}{2}-RL,\\
 \kappa_{0,2}&=-L+R+\frac{R^2}{2}-\frac32.
\end{align*}
At $\lambda=1$ the contact is a polynomial of degree at most $r-1$
in $u$. Consequently $\kappa_{m,r}(1,\mu)=0$ whenever $m\ge r$.
At $\lambda=\mu=1$ all contacts vanish. These are exact truncations
of the correction, not asymptotic statements about its size.

\subsection{The complete singular power--logarithm expansion}

The following evaluated germ specifies every divergent term as well
as the constant. Put
\begin{equation}\label{urc:local-decomposition}
 P_r(u)=(-1)^r(1+u)_r,\qquad
 Z_r(u,x)=\partial_x^r Z(u,x)
         =P_r(u)x^{-r-1-u}+A_r^{\mathrm{reg}}(u,x),
\end{equation}
where $A_r^{\mathrm{reg}}(u,x)=\partial_x^r Z(u,1+x)$.
The superscript distinguishes this analytic function from the finite
polynomial $A_r$ above. Define
\begin{equation}\label{urc:analytic-coefficients}
 a_{t,k}(v)=\frac1{k!}\partial_z^{t+k}Z(v,z)\big|_{z=1}.
\end{equation}
These are holomorphic near $v=0$, and
\[
 (-1)^n\partial_v^n a_{t,k}(0)
      =\frac{\gamma_n^{(t+k)}(1)}{k!}.
\]

The only nonpolynomial local kernel needed is the elementary beta
integral
\begin{align}
 B_{r,t}(u,v;\lambda,\mu)
 &=P_r(u)P_t(v)\int_0^\infty
 \frac{dy}{(y+\lambda)^{r+1+u}(y+\mu)^{t+1+v}}\notag\\
 &=\frac{P_r(u)P_t(v)\mu^{-r-t-1-u-v}}{r+t+1+u+v}\notag\\
 &\quad\cdot{}_2F_1\left(
 \begin{matrix}r+1+u,\ r+t+1+u+v\\r+t+2+u+v\end{matrix};
                   1-\frac\lambda\mu\right).
 \label{urc:beta-kernel}
\end{align}
The integral fixes its branch. It is holomorphic on a neighborhood
of $(u,v)=(0,0)$, including $\lambda=\mu$.

For $0\le k<r$, define the finite expression
\begin{equation}\label{urc:mixed-kernel}
 M_{r,k}(u;\lambda,\mu)
 =P_r(u)\lambda^{k-r-u}
  \sum_{j=0}^k\binom kj
      \frac{(\mu/\lambda-1)^j}{r-k+j+u}.
\end{equation}
For the critical degree use the removable germ
\begin{equation}\label{urc:critical-kernel}
 N_r(u;\epsilon,\lambda,\mu)
 =P_r(u)\left\{
 \epsilon^{-u}\lambda^{-u}
 \left(\frac1u+\sum_{j=1}^r\binom rj
                  \frac{(\mu/\lambda-1)^j}{u+j}\right)
 -\frac1u\right\}.
\end{equation}
At each fixed $\epsilon>0$ its apparent pole at $u=0$ cancels
before any differentiation. Its constant in $\log\epsilon$ is
$r!C_r(u;\lambda,\mu)$.

Finally let $\mathcal F_{r,t}(u,v)$ denote the holomorphic germ whose
coefficient dictionary is
\begin{equation}\label{urc:fp-germ}
 (-1)^{m+n}\partial_u^m\partial_v^n
          \mathcal F_{r,t}(0,0)=F_{m,n}^{r,t}.
\end{equation}
To specify it independently of a meromorphic regularization, set
\begin{align}
 R_{r,t}(u,v;x)={}&Z_r(u,x)Z_t(v,x)
 -P_r(u)P_t(v)x^{-r-t-2-u-v}\notag\\
 &-P_r(u)\sum_{k=0}^r a_{t,k}(v)x^{k-r-1-u}
 -P_t(v)\sum_{k=0}^t a_{r,k}(u)x^{k-t-1-v},
 \label{urc:remainder}
\end{align}
and define
\begin{align}
 \mathcal F_{r,t}(u,v)={}&\int_0^1R_{r,t}(u,v;x)\,dx
 -\frac{P_r(u)P_t(v)}{r+t+1+u+v}\notag\\
 &-P_r(u)\sum_{k=0}^{r-1}\frac{a_{t,k}(v)}{r-k+u}
 -P_t(v)\sum_{k=0}^{t-1}\frac{a_{r,k}(u)}{t-k+v}.
 \label{urc:fp-germ-construction}
\end{align}
Empty sums are zero. The critical terms $k=r$ and $k=t$ have no
added constant: every monomial $x^{-1}\log^h x$ has finite part
zero on $(0,1)$ in the coordinate \eqref{urc:finitepart}. This is
why retaining a whole meromorphic $1/u$ expression would give the
wrong Stieltjes coefficient convention.

\begin{theorem}[All singular terms in one joint germ]
\label{urc:full-theorem}
For fixed $r,t\ge0$ and $\lambda,\mu>0$, as $\epsilon\downarrow0$,
\begin{align}
 &\int_0^1Z_r(u,x+\lambda\epsilon)Z_t(v,x+\mu\epsilon)\,dx
 =\mathcal F_{r,t}(u,v)
   +\epsilon^{-r-t-1-u-v}B_{r,t}(u,v;\lambda,\mu)\notag\\
 &\quad+\sum_{k=0}^{r-1}a_{t,k}(v)
               \epsilon^{k-r-u}M_{r,k}(u;\lambda,\mu)
       +\sum_{k=0}^{t-1}a_{r,k}(u)
               \epsilon^{k-t-v}M_{t,k}(v;\mu,\lambda)\notag\\
 &\quad+a_{t,r}(v)N_r(u;\epsilon,\lambda,\mu)
       +a_{r,t}(u)N_t(v;\epsilon,\mu,\lambda)+o(1).
 \label{urc:full-expansion}
\end{align}
The remainder tends to zero locally uniformly for $(u,v)$ in a
complex neighborhood of zero. Every fixed mixed order derivative
may therefore be taken in this asymptotic formula. Applying
$(-1)^{m+n}\partial_u^m\partial_v^n$ and setting $u=v=0$ gives
the complete negative-power and logarithmic expansion of
\eqref{urc:integral}, through its constant.
\end{theorem}

The possible negative powers are especially sparse. They consist of
$\epsilon^{-(r+t+1)}$, together with $\epsilon^{-r},\ldots,
\epsilon^{-1}$ and $\epsilon^{-t},\ldots,\epsilon^{-1}$ when the
corresponding derivative orders are positive. The first power has
a logarithmic polynomial of degree at most $m+n$. The two finite
mixed families have degrees at most $m$ and $n$. The critical
terms give logarithmic polynomials of degrees at most $m+1$ and
$n+1$. Formula \eqref{urc:full-expansion} includes cancellations
between these polynomials automatically.

\begin{proof}
The Hurwitz recurrence gives \eqref{urc:local-decomposition} exactly.
Its analytic part has the Taylor expansion
\[
 A_t^{\mathrm{reg}}(v,z)=\sum_{k=0}^r a_{t,k}(v)z^k
                              +O(z^{r+1})
\]
uniformly for $v$ in a sufficiently small closed complex disk.
Expand the product of its singular and analytic parts. For the
singular--singular term, substitute $x=\epsilon y$ and complete
the interval to $(0,\infty)$. Its omitted tail tends locally
uniformly to $P_r(u)P_t(v)/(r+t+1+u+v)$, giving precisely the
beta term in \eqref{urc:full-expansion} and the first rational
constant in \eqref{urc:fp-germ-construction}.

For a mixed Taylor monomial the finite identity is
\begin{align}
 &P_r(u)\int_0^1
       (x+\lambda\epsilon)^{-r-1-u}(x+\mu\epsilon)^k\,dx
 \notag\\
 &\quad=P_r(u)\sum_{j=0}^k\binom kj
       ((\mu-\lambda)\epsilon)^j
 \frac{(\lambda\epsilon)^{k-j-r-u}
                  -(1+\lambda\epsilon)^{k-j-r-u}}
      {r-k+j+u}.
 \label{urc:exact-mixed-primitive}
\end{align}
It follows by expanding $x+\mu\epsilon=(x+\lambda\epsilon)
+(\mu-\lambda)\epsilon$ and integrating each power. At $k<r$,
the lower endpoint is $\epsilon^{k-r-u}M_{r,k}$, while the upper
endpoint tends to $-P_r(u)/(r-k+u)$. At $k=r$ the combination of
the lower endpoint and $-P_r(u)/u$ is exactly $N_r$.
The remaining upper-endpoint error is holomorphic at $u=0$ and
tends uniformly to zero; in particular,
$\{1-(1+\lambda\epsilon)^{-u}\}/u=O(\epsilon)$ uniformly near
zero. Exchange the two factors for the other mixed terms.

It remains to justify passage to the limit in the Taylor remainders.
For $|\Re u|,|\Re v|\le\eta<1/4$ and fixed positive rates, a
singular factor times the Taylor remainder is bounded by
$C(x+\epsilon)^{-\eta}$, hence by $Cx^{-\eta}$ on $(0,1)$.
This integrable majorant is uniform on compact parameter disks.
The product of analytic remainders is uniformly bounded. Dominated
convergence therefore gives the integral in
\eqref{urc:fp-germ-construction}, locally uniformly in $u,v$.
Cauchy's formula on slightly larger parameter disks gives the
same conclusion after every prescribed mixed order derivative.
This proves \eqref{urc:full-expansion}, including its joint rather
than merely raywise spectral meaning.

Coefficient extraction in \eqref{urc:fp-germ-construction} gives
\eqref{urc:fp-germ}: noncritical singular monomials have their
ordinary finite primitives at one, whereas every critical
$x^{-1}\log^h x$ contributes zero. Thus no exchange between an
unsubtracted divergent integral and a spectral derivative is used.
\end{proof}

\begin{proof}[Proof of Theorem~\ref{urc:constant-theorem}]
In \eqref{urc:full-expansion}, the beta term and the noncritical
mixed terms have strictly negative powers of $\epsilon$, also
after spectral differentiation. They cannot contribute to the
constant under discussion. In each critical term, retain the
constant in $\log\epsilon$, namely $r!C_r$ or $t!C_t$.
Equation \eqref{urc:analytic-coefficients} then gives
\eqref{urc:constant-formula}. At $u=0$,
\[
 C_r(0;\lambda,\mu)=(-1)^r\left\{-\log\lambda+
       \sum_{j=1}^r\binom rj\frac{(R-1)^j}{j}\right\}.
\]
The polynomial identity
\[
 \sum_{j=1}^r\binom rj\frac{(R-1)^j}{j}
       =\sum_{j=1}^r\frac{R^j-1}{j}
\]
follows by differentiating both sides in $R$ and using their common
value zero at $R=1$. This proves \eqref{urc:zeroth-contact}.
At $r=0$, expanding $(\lambda^{-u}-1)/u$ gives
\eqref{urc:no-argument-derivatives}.
\end{proof}

\subsection{Polylogarithmic coefficients of the divergent beta term}

The hypergeometric notation in \eqref{urc:beta-kernel} does not add
an unspecified family of constants. It is enough to differentiate
the base kernel in the shifts:
\[
 B_{r,t}(u,v;\lambda,\mu)
       =\partial_\lambda^r\partial_\mu^t
                      B_{0,0}(u,v;\lambda,\mu).
\]
For $\lambda>\mu$, put $d=\lambda-\mu$ and $z=d/\lambda$.
The substitution $y=d/(x+\lambda)$ gives
\begin{equation}\label{urc:polylog-beta}
 B_{0,0}(u,v;\lambda,\mu)
  =d^{-1-u-v}\int_0^z
       \frac{\exp\{(u+v)\log y-v\log(1-y)\}}{1-y}\,dy.
\end{equation}
Every coefficient of total spectral degree $N$ is a finite linear
combination of harmonic polylogarithms on the alphabet $\{0,1\}$
and logarithms, of weight at most $N+1$. This follows by shuffling
the logarithmic factors and integrating once with $dy/(1-y)$.
Shift differentiation preserves this finite function class. The
case $\mu>\lambda$ follows by exchange; the coincident case is
the elementary expression
\[
 B_{r,t}(u,v;\lambda,\lambda)
 =\frac{P_r(u)P_t(v)\lambda^{-r-t-1-u-v}}{r+t+1+u+v}.
\]
Thus every coefficient in the full subtraction belongs to an
explicit finite alphabet of harmonic numbers, logarithms,
polylogarithms, and Hurwitz or Stieltjes derivatives at one.

\subsection{Three completely explicit polygamma collision laws}

Write
\begin{equation}\label{urc:D}
 D(\lambda,\mu)=\frac{\log(\lambda/\mu)}{\lambda-\mu},
 \qquad D(\lambda,\lambda)=\frac1\lambda,
\end{equation}
and let $D_\lambda$ and $D_{\lambda\mu}$ be its indicated
derivatives. Explicitly, away from the diagonal,
\begin{align*}
 D_\lambda&=\frac{(\lambda-\mu)/\lambda-\log(\lambda/\mu)}
                  {(\lambda-\mu)^2},\\
 D_{\lambda\mu}
 &=\frac{(\lambda-\mu)(\lambda^{-1}+\mu^{-1})
                       -2\log(\lambda/\mu)}{(\lambda-\mu)^3}.
\end{align*}
Their diagonal values are $-1/(2\lambda^2)$ and
$1/(3\lambda^3)$. The inherited coincident evaluations \cite{LatticeIncoming} are
\[
 F_{0,0}^{0,0}=2\gamma_1-2\zeta(2),\qquad
 F_{0,0}^{1,0}=\zeta(2),\qquad
 F_{0,0}^{1,1}=2\zeta(3)+4\zeta'(3).
\]
Using these established values in the new collision formula gives
\begin{align}
 \int_0^1\psi(x+\lambda\epsilon)\psi(x+\mu\epsilon)\,dx
 ={}&\frac{D(\lambda,\mu)}\epsilon-2\gamma\log\epsilon
       +2\gamma_1-2\zeta(2)-\gamma\log(\lambda\mu)+o(1),
 \label{urc:digamma-example}\\
 \int_0^1\psi'(x+\lambda\epsilon)\psi(x+\mu\epsilon)\,dx
 ={}&\frac{D_\lambda(\lambda,\mu)}{\epsilon^2}
       -\frac{\gamma}{\lambda\epsilon}
       +\boxed{\zeta(2)\left\{\frac\mu\lambda+
                               \log\frac\mu\lambda\right\}}
       +o(1),
 \label{urc:ordered-example}\\
 \int_0^1\psi'(x+\lambda\epsilon)\psi'(x+\mu\epsilon)\,dx
 ={}&\frac{D_{\lambda\mu}(\lambda,\mu)}{\epsilon^3}
       +\frac{\zeta(2)}\epsilon
                       \left(\frac1\lambda+\frac1\mu\right)
       +4\zeta(3)\log\epsilon\notag\\
 &+4\zeta'(3)+2\zeta(3)
       \left\{3+\log(\lambda\mu)-\frac\mu\lambda-
                                      \frac\lambda\mu\right\}+o(1).
 \label{urc:trigamma-example}
\end{align}
The ordered law \eqref{urc:ordered-example} is invariant under a
common scaling of the two rates. Its exchange companion has
constant $\zeta(2)\{\lambda/\mu+\log(\lambda/\mu)\}$. Their
sum equals $\zeta(2)(\lambda/\mu+\mu/\lambda)$, in exact
agreement with the endpoint derivative of the digamma product.
This independently checks the logarithmic signs and the ratio
terms; exchange symmetry alone would not determine either ordered
constant separately.

\subsection{A first Stieltjes collision with a dilogarithm}

Let $\gamma_j=\gamma_j(1)$ and put
\begin{equation}\label{urc:E}
 E(\lambda,\mu)=
 \frac{\log(\lambda/\mu)\log\lambda+
               \operatorname{Li}_2(1-\mu/\lambda)}{\lambda-\mu},
 \qquad E(\lambda,\lambda)=\frac{1+\log\lambda}{\lambda}.
\end{equation}
This is the ordinary convergent integral of
$\log(y+\lambda)/\{(y+\lambda)(y+\mu)\}$ on $(0,\infty)$,
as follows either from \eqref{urc:polylog-beta} or one integration
by parts. The standard coincident product germ gives
\begin{equation}\label{urc:F10}
 F_{1,0}^{0,0}=\frac32\gamma_2+2\gamma\zeta(2)-\zeta(3).
\end{equation}
Consequently the complete first-index expansion is
\begin{align}
 \int_0^1\gamma_1(x+\lambda\epsilon)
                    \gamma_0(x+\mu\epsilon)\,dx
 ={}&\frac{D(\lambda,\mu)\log\epsilon+E(\lambda,\mu)}\epsilon
       -\frac\gamma2(\log\epsilon)^2\notag\\
 &-(\gamma\log\lambda+\gamma_1)\log\epsilon
       +\frac32\gamma_2+2\gamma\zeta(2)-\zeta(3)\notag\\
 &-\frac\gamma2(\log\lambda)^2-\gamma_1\log\mu+o(1).
 \label{urc:first-Stieltjes-example}
\end{align}
All logarithms and the dilogarithm in this formula are real for
positive rates, with the displayed continuous diagonal limit.

For completeness, \eqref{urc:F10} can be extracted without a
numerical relation search. Put $w=u+v$,
$Q=\Gamma(1+w)/\{\Gamma(1+u)\Gamma(1+v)\}$ and
$E_0(u)=\pi u/\sin\pi u$. The inherited coincident Hurwitz
product formula gives the holomorphic finite-part germ
\begin{align*}
 \mathcal F_{0,0}(u,v)
 ={}&\frac{w-Q\{uE_0(v)+vE_0(u)\}}{uvw}\\
 &+\frac{Z(v,1)-QE_0(u)Z(w,1)}u
  +\frac{Z(u,1)-QE_0(v)Z(w,1)}v.
\end{align*}
Each divided difference is removable. Its degree-zero term is
$2\gamma_1-2\zeta(2)$ and its coefficient of $u$ is
$\zeta(3)-3\gamma_2/2-2\gamma\zeta(2)$. Applying the sign
convention \eqref{urc:stieltjes} proves \eqref{urc:F10}.

\subsection{Verification, scope, and further questions}

The accompanying verifier checks the finite mixed primitive,
the rising-factorial cancellations, the compact ratio polynomial,
and contact jets by two exact coefficient constructions. Separate
high-precision ordinary quadratures check the three beta kernels
and the three polygamma collisions at diagonal, unequal integral,
and unequal rational rates. The numerical collision residuals are
compared at decreasing positive $\epsilon$ after the complete
displayed subtraction. They test the asymptotic constants and
normalization; they are not certified error enclosures.

The result concerns fixed positive rates. Paths with a vanishing or
unbounded ratio require an additional asymptotic regime; inserting
a rate that depends on $\epsilon$ into a fixed-compactness proof
would not justify that regime. The following are concrete next
questions.
\begin{enumerate}
\item Determine the complete transition for
$a=\lambda\epsilon^\alpha$, $b=\mu\epsilon^\beta$ with
$\alpha\ne\beta$, including the boundary between the two
scales. The explicit finite ratio polynomial and the beta kernel
make every potential extra power visible.
\item Extend the locality theorem to three independently shifted
Stieltjes factors. With more than two factors, products of two
singular parts and one analytic Taylor coefficient can already
affect the constant; the required kernel should be stated before
any reduction to a smaller alphabet is claimed.
\item Classify linear combinations of ordered derivative products
whose collision constant depends only on the ratio $\mu/\lambda$,
or is independent of the rates. Formula \eqref{urc:finite-coefficients}
reduces this to an exact finite polynomial problem at each index.
\item Continue the rates to a common complex sector and derive
the jumps when the lower-endpoint logarithms cross their chosen
cuts. The local beta integral supplies a natural branch convention.
\end{enumerate}

% END sections/04_collisions.tex

% BEGIN sections/05_audit_and_questions.tex
\section{Source corrections, integration decisions, and further research}
\label{audit:section}

\subsection{Precise changes to the source status}

No newly false inherited mathematical theorem was found in the
proof dependencies used here. The proposed changes are extensions
of valid theorems and corrections to research-status lists.
They should be applied with their original attribution intact.

\begin{enumerate}
\item In the ordered-boundary report's Section~9, replace the
statement that complete gap-six coverage is open by a reference
to Theorem~\ref{gs:main}, explicitly recording its
computation-assisted proof. Retain the broader all-pairs question.
The source's stronger sign conjecture
$f_b>0$ for every $b\ge7$ is not proved by a nonzero finite-field
residue. We establish positivity for $b\ge2048$ and nonvanishing
at all earlier admissible indices; we make no additional finite
sign claim.
\item In the bilinear determinant report's Section~10, mark the
complex rational-weight degree--rank question solved by
Theorem~\ref{cr:degree}. Broaden the generic rational-weight
degree statement to complex coefficients. Preserve the real
eventual-sign hypotheses in the source's sharper numerical
exception bound; that bound does not follow from the new proof.
\item In the lattice-contact report's Question~Q9 and the final
paragraph of its independent-shift section, mark the fixed
positive unequal-rate collision question solved by
Theorems~\ref{urc:constant-theorem} and \ref{urc:full-theorem}.
Keep nonlinear paths and ratios tending to zero or infinity open.
\item Do not reopen the already proved smaller-order-ten and
smaller-order-twelve cases. Do not relabel the arbitrary mixed
even-tail product calculus or the existing Stieltjes spectral-face
results as new consequences of this article.
\end{enumerate}

Several implementation conventions deserve explicit retention.
In the late Bernoulli coefficient matrix, a lower binomial index
of zero is omitted, rather than evaluated as $\binom n0=1$.
Changing that convention changes the theorem's matrix. In the
complex rank proof the relevant bilinear operation is ordinary
transposition; replacing it by conjugate transposition would
change the rational algebraic problem. In the Stieltjes section
the constant depends on the stated small parameter; critical
Laurent quotients must be removed jointly before taking spectral
derivatives. These are guards against incorrect integrations of
the results, not allegations that the cited source theorems made
those mistakes.

The manuscript's old rejected $S_8$ candidate should remain
rejected, and its revised $S_8$ and short $S_6$ candidates should
retain their conjectural status. Nothing in the present rank,
Hahn, or collision arguments implies those period equalities.
A small numerical residual does not supply their missing proof.

\subsection{A proposed research program}

The next questions below are selected for their potential to
produce exact identities or structural theorems. Each has a
concrete starting point supplied by the article.

\begin{enumerate}[label=\textbf{Q\arabic*.},leftmargin=*]
\item \textbf{Complete the gap-eight family.}
For half-gap $m=4$, make the inherited four-mode Gamma basis
effective. Two modes have a vanishing leading imaginary phase.
Derive a uniform two-term estimate for each, control the
four-by-four limiting determinant, and obtain an explicit finite
completion threshold. The target is nonvanishing for every
$b\ge5$, not simply a larger numerical table.

\item \textbf{Replace the gap-six finite completion by a structural proof.}
Find a recurrence invariant or a variation-diminishing theorem
which proves $f_b\ne0$ at every admissible $b$. A fixed-sign
assertion from the first index cannot work because the initial
sign transition is real. The stronger source conjecture
$f_b>0$ for $b\ge7$ would also settle the sign issue left
unresolved by the modular certificate.

\item \textbf{Connect fixed-gap and fixed-smaller-order regimes.}
The source has eventual theorems in both directions, and the
present proof makes one gap ray effective. Develop a
two-parameter determinant expansion uniform when both indices
increase. The limiting reciprocal-square nodes become less
separated as the gap grows; merely reusing fixed-gap constants
does not cover the intervening region.

\item \textbf{Evaluate the Gamma normalization determinants.}
The initial determinant is nonzero by an integral alternant
argument. Seek explicit Gamma, hypergeometric, or recurrence
evaluations for its normalization $K_m$ at $m=2,3,\ldots$.
Such evaluations would turn the recurrence certificates into
more compact identities rather than just nonvanishing tests.

\item \textbf{Classify rationally summable polynomial subspaces.}
For a given rational weight $R$, classify the spaces $V$ on
which every discrete-residue bilinear form of $Rfg$ vanishes.
Theorem~\ref{cr:primitive} proves nondegeneracy of every rational
primitive once $V$ is known; it does not classify these spaces.
This asks for a concrete bridge between rational summability
and symmetric pencils with varying kernels.

\item \textbf{Find intrinsic complex cutoff-exception invariants.}
Express the degree or factorization of the exceptional-cutoff
polynomial directly from the pole orders, residues, and integer
dispersions of $R$, without constructing the full determinant.
The coefficient-gcd criterion is exact, but a residue-level
description could explain and predict the later nonreal example
in Section~\ref{cr:section}.

\item \textbf{Extend the adjacent-pole compiler to several poles.}
For $R(n)=\sum_{\ell=1}^d c_\ell/(n+a_\ell)$, determine when
the highest harmonic coefficient reduces to a determinant of
fixed size as $r$ varies. For integer-separated poles the
compressed measure is a polynomial transform, suggesting a
Hahn or multiple-orthogonal representation. The difficult
part is to classify endpoint pole cancellations and exceptional
residue varieties, rather than invoke the general Christoffel
formula alone.

\item \textbf{Classify when classical polylogarithms suffice.}
Theorem~\ref{hp:all-r-generator} gives an entire arbitrary-rank
family using only one dilogarithm. Determine necessary and
sufficient conditions on a rational weight for its cutoff
generators to remain in a fixed classical-polylogarithm space.
Generic harmonic degree supplies an upper-level clue but is
not by itself a minimal-depth theorem.

\item \textbf{Evaluate the Euler primitives of the ensemble family.}
For the explicit $\mathcal Z_N$, determine all endpoint values
$\sum_{M\ge1}\mathcal Z_{M-1}/M^s$ in their convergence range,
and their alternating companions. Equation
\eqref{hp:euler-reduction} generalizes to a small list of
linear and quadratic Euler sums. The question is when these
reduce to ordinary zeta products and when irreducible higher
weight components first survive.

\item \textbf{Resolve genuinely different collision scales.}
Derive the full expansion for
$a=\lambda\epsilon^\alpha$ and $b=\mu\epsilon^\beta$,
$\alpha\ne\beta$. The fixed-ratio theorem cannot simply be
used with an $\epsilon$-dependent ratio without new uniform
analysis. Its beta kernel and finite contact polynomial
identify the terms which must enter a transition formula.

\item \textbf{Develop three-factor collision constants.}
For three independently shifted Stieltjes derivatives,
products of two singular factors with one analytic coefficient
may affect the constant. Determine the finite local contact
tensor and the appropriate multivariable beta kernel before
claiming a reduction to the two-factor alphabet.

\item \textbf{Identify invariant ordered combinations and complex jumps.}
Classify linear combinations of derivative products whose
collision constants depend only on the rate ratio, or are
independent of both rates. Then continue the rates in a
common complex sector and compute logarithmic jumps across
the boundary cuts. At fixed orders the first task is an
explicit finite polynomial problem in the contact coefficients.
\end{enumerate}

The strongest unresolved global targets remain the unrestricted
moment-prefix conjecture and the manuscript's distinguished
Gaussian period identities. The present article supplies three
completed local research programs and concrete next steps toward
broader statements; it does not infer those global targets from
finite evidence.


% END sections/05_audit_and_questions.tex

\appendix
% BEGIN sections/05a_gamma_appendix.tex
% Specialized proof of the inherited analytic input for the gap-six theorem.
% All labels use ga:. This is an attribution-preserving specialization, not a new claim.
\section{The Gamma basis used in the gap-six proof}
\label{ga:appendix}

This appendix reproduces the case $m=3$ of the Gamma-basis theorem
in the incoming \emph{Ordered Boundary Identities and Gamma Kernels}
report~\cite[Section~8]{OrderedIncoming}.  Its purpose is to make the
analytic input to Theorem~\ref{gs:main} independently reproducible
within this article.  The contour construction and the initial
determinant argument are inherited results; the effective estimates
and complete gap-six conclusion are established in Section~\ref{gs:section}.

Put
\[
 W(x)=\left|\Gamma\!\left(\tfrac14+\tfrac{ix}{2}\right)\right|^{12},
 \quad P_1(x)=\sinh(2\pi x),\quad P_2(x)=\sinh(\pi x),
 \quad P_3(x)=\tanh(\pi x).
\]
For $1\le j\le3$, define
\begin{equation}\label{ga:G}
 G_j(z)=4\sin(z/2)\int_0^\infty\sinh(xz)P_j(x)W(x)\,dx,
 \qquad t_{n,j}=[z^{2n}]G_j(z).
\end{equation}
The density estimate of Lemma~\ref{gs:density}, or the ordinary
Gamma modulus asymptotic~\cite[Equation~5.11.9]{DLMF}, proves
locally uniform convergence in $|\operatorname{Re}z|<j\pi$.
These functions are even and vanish at zero.

\begin{proposition}[Specialized Gamma basis]\label{ga:basis}
The three germs in \eqref{ga:G} satisfy
\begin{equation}\label{ga:null}
 \bigl(G_j^{(6)}-G_j^{(6)}(0)\bigr)\cot(z/2)
                    +(G_j\cot(z/2))^{(6)}=0.
\end{equation}
Their sixth derivatives at zero are
\begin{equation}\label{ga:top}
 G_1^{(6)}(0)=G_2^{(6)}(0)=0,
 \qquad G_3^{(6)}(0)=256\pi^3.
\end{equation}
The matrix $T=(t_{n,j})_{1\le n,j\le3}$ has negative determinant.
Thus these germs form a basis of the formal null space of
\eqref{ga:null}, and their later coefficient rows satisfy
\[
                     (t_{b+i,j})_{i,j=1}^3=F_bT
\]
with the rational fundamental matrix of Section~\ref{gs:section}.
\end{proposition}

\subsection{The Fourier shift and its two residues}

For one of the three odd multipliers $P$, set
\[
 \mathcal F_P(t)=\int_{\mathbb R}e^{ixt}P(x)W(x)\,dx,
 \qquad \mathcal Q_P(t)=2\sinh(t/2)\mathcal F_P(t),
 \qquad D=\frac{d}{dt}.
\]
The integrals and all derivatives converge absolutely for real $t$
and define analytic germs.  The oddness of $P$ makes $\mathcal F_P$
odd, so $\mathcal Q_P$ is even and vanishes at zero.  Define
\[
 \mathcal E_P(t)=e^t(D+\tfrac12)^6\mathcal F_P(t)
                        -e^{-t}(D-\tfrac12)^6\mathcal F_P(t).
\]
Since
\[
 \mathcal Q_P=(e^{t/2}-e^{-t/2})\mathcal F_P,
 \qquad \mathcal Q_P\coth(t/2)
              =(e^{t/2}+e^{-t/2})\mathcal F_P,
\]
ordinary differentiation gives the exact identity
\begin{equation}\label{ga:operator}
 \coth(t/2)\mathcal Q_P^{(6)}(t)
       +(\mathcal Q_P(t)\coth(t/2))^{(6)}
                  =\frac{\mathcal E_P(t)}{\sinh(t/2)}.
\end{equation}

For complex arguments, use the single-valued meromorphic continuations
\begin{align*}
 W(y)&=\left[\Gamma(\tfrac14+iy/2)
                       \Gamma(\tfrac14-iy/2)\right]^6,\\
 V(y)&=\left[\Gamma(\tfrac34+iy/2)
                       \Gamma(\tfrac34-iy/2)\right]^6.
\end{align*}
The absolute-value notation for $W$ is used only on the real axis.
Gamma recurrence gives
\begin{equation}\label{ga:recurrence}
 (iy-\tfrac12)^6W(y+i)
   =(iy+\tfrac12)^6W(y-i)=64V(y).
\end{equation}
In the first Fourier integral defining $\mathcal E_P$, substitute
$y=x-i$, so its contour is $\mathbb R-i$; in the second substitute
$y=x+i$, so its contour is $\mathbb R+i$.  These substitutions
remove the respective $e^{\pm t}$ factors.  Moreover
$P(y+i)=P(y-i)$ for every allowed $P$.  Thus bottom minus top
is the integral of the same meromorphic function
\begin{equation}\label{ga:common-integrand}
                  64e^{iyt}V(y)P(y+i)
\end{equation}
on the two horizontal sides of a positively oriented rectangle.

The nearest poles of $V$ are $\pm3i/2$, outside the closed strip
$|\operatorname{Im}y|\le1$.  The apparent closer Gamma poles in
the shifted $W$ expressions have already been canceled exactly
by \eqref{ga:recurrence}.  For a sinh multiplier, $P(y+i)$ is
an entire function: it equals $\sinh(2\pi y)$ for $P_1$ and
$-\sinh(\pi y)$ for $P_2$.  For $P_3$ it equals
$\tanh(\pi y)$, with precisely two poles in the strip,
$y=\pm i/2$.
Uniform Stirling estimates on the closed strip bound the integrand
by a polynomial times $e^{-(3-k)\pi|\operatorname{Re}y|}$ in
the sinh case, where $k=2$ or $1$, and by a polynomial times
$e^{-3\pi|\operatorname{Re}y|}$ in the tanh case.
For fixed real $t$, the factor $e^{iyt}$ is bounded there.
The vertical sides therefore vanish as their real parts tend to
infinity.  All contour shifts are consequently justified.

There are no enclosed poles for $P_1,P_2$, so
$\mathcal E_{P_1}=\mathcal E_{P_2}=0$.
At each tanh pole the residue is $1/\pi$, and
\[
                         V(i/2)=V(-i/2)=\pi^3.
\]
The residues of \eqref{ga:common-integrand} are therefore
$64\pi^2e^{-t/2}$ and $64\pi^2e^{t/2}$.  Bottom minus top is
$2\pi i$ times their sum, giving
\begin{equation}\label{ga:residue}
                \mathcal E_{P_3}(t)=256i\pi^3\cosh(t/2).
\end{equation}

The second term on the left of \eqref{ga:operator} is holomorphic
at zero, since $\mathcal Q_P(t)=O(t^2)$.  The first has Laurent
coefficient $2\mathcal Q_P^{(6)}(0)$ at $t^{-1}$.
Comparing \eqref{ga:operator} with \eqref{ga:residue} yields
\[
 \mathcal Q_{P_1}^{(6)}(0)=\mathcal Q_{P_2}^{(6)}(0)=0,
 \qquad\mathcal Q_{P_3}^{(6)}(0)=256i\pi^3.
\]
The exact identities also show that
\[
 (\mathcal Q_P^{(6)}-\mathcal Q_P^{(6)}(0))\coth(t/2)
                   +(\mathcal Q_P\coth(t/2))^{(6)}=0.
\]
Pairing positive and negative $x$ gives
$G_j(z)=i\mathcal Q_{P_j}(iz)$.  Rotation uses
$\coth(iz/2)=-i\cot(z/2)$, proving \eqref{ga:null}.
Since $i\cdot i^6=-i$, it also gives \eqref{ga:top}.

\subsection{The initial determinant has an exact positive integral}

Write $\phi_n(x)=\operatorname{Im}(x+i/2)^{2n}$.  The elementary
identity
\[
 \sin(z/2)\sinh(xz)=\sum_{n\ge1}\frac{\phi_n(x)}{(2n)!}z^{2n}
\]
and local uniform domination imply
\[
               t_{n,j}=\frac4{(2n)!}
                   \int_0^\infty\phi_n(x)P_j(x)W(x)\,dx.
\]
The polynomial $\phi_n$ has leading term $nx^{2n-1}$ and only
lower odd powers.  A triangular row transformation consequently gives
\[
 \det T=\left(\prod_{n=1}^3\frac{4n}{(2n)!}\right)
   \det\left(\int_0^\infty x^{2n-1}P_j(x)W(x)\,dx\right)_{n,j=1}^3,
 \qquad\prod_{n=1}^3\frac{4n}{(2n)!}=\frac1{90}.
\]
For $\Delta(y_1,y_2,y_3)=\prod_{i<j}(y_j-y_i)$,
continuous Cauchy--Binet gives the following exact normalization:
\begin{equation}\label{ga:initial-determinant}
 \boxed{\begin{aligned}
 \det T={}&-\frac1{270}\int_{(0,\infty)^3}
       \Delta(x_1^2,x_2^2,x_3^2)\\
 &\qquad\cdot\Delta(\cosh\pi x_1,\cosh\pi x_2,\cosh\pi x_3)
       \prod_{r=1}^3x_rW(x_r)\tanh(\pi x_r)\,d\boldsymbol{x}<0.
 \end{aligned}}
\end{equation}
Indeed $\det(x_r^{2n-1})=(\prod_rx_r)\Delta(x_1^2,x_2^2,x_3^2)$.
After dividing row $r$ of $\det(P_j(x_r))$ by
$\tanh(\pi x_r)$ and putting $v_r=\cosh(\pi x_r)$, its
columns are $2v_r^2,v_r,1$.  Its determinant is
$-2\Delta(v_1,v_2,v_3)$; together with the Cauchy--Binet factor
$1/3!$ and $1/90$ this gives $-1/270$.
The integral converges absolutely by exponential decay.  Both
Vandermondes have the same sign whenever the $x_r$ are distinct,
since $x\mapsto x^2$ and $x\mapsto\cosh(\pi x)$ are strictly
increasing on $(0,\infty)$.  Its integrand is therefore positive
off the collision planes, proving strict negativity.

The coefficient recurrence of Section~\ref{gs:section} gives a
three-dimensional formal null space with arbitrary first three
coefficients.  Equation \eqref{ga:initial-determinant} proves that
our three germs span it, and coefficient comparison proves
$(t_{b+i,j})=F_bT$, completing Proposition~\ref{ga:basis}.
As a normalization check, \eqref{ga:top} also gives the inherited
ordinary Gamma moments
\[
 \int_0^\infty\left(3x^5-\frac52x^3+\frac3{16}x\right)
                 P_j(x)W(x)\,dx
 =\begin{cases}0,&j=1,2,\\64\pi^3,&j=3.\end{cases}
\]

% END sections/05a_gamma_appendix.tex

% BEGIN sections/06_worked_example.tex
% Exact ordinary convergent example. K is encoded by its integer coefficient row.
\section{An exact \texorpdfstring{$(14,20)$}{(14,20)} polygamma-tail evaluation}
\label{example:section}
For $(p,q)=(14,20)$ put $K(y)=y^{21}\sum_{j=0}^{9}c_jy^{2j}$,
where the coefficient row $(c_0,\ldots,c_9)$ is
\[\begin{aligned}
-442272850709916359100&,\quad -916343865300090386578\\
7827122405129750301156&,\quad -17591879878937090990204\\
24015626210965489858160&,\quad -27581453408920271748009\\
51508173856096311750572&,\quad 53264037121751700169882\\
9020688685919084220372&,\quad 346718746623385922565
\end{aligned}\]
Then the centered difference $P(x)=K(x+1/2)-K(x-1/2)$ and the polynomial
\[\begin{aligned}
 S(x)={}&\frac{1040156239870157767695}{19}x^{6}
       +\frac{1618580484745959817903323}{494}x^{4}\\
       &+\frac{171589487161427755057692583}{3952}x^{2}
       +\frac{25964774156256502396231852}{247}.
\end{aligned}\]
give the exact absolutely convergent identity
\[\begin{aligned}
 \sum_{n\ge0}\left\{\frac{P(n+1/2)}{13!\,19!}
 \psi^{(13)}(n+1)\psi^{(19)}(n+1)-S(n+1/2)\right\}\\
 =-\frac{70897571597555316074547}{2}.
\end{aligned}\]
The complete finite telescoper is recorded in \texttt{gap\_six\_14\_20\_example.json} and regenerated by \texttt{create\_gap\_six\_example.py}. Its verification uses only polynomial and rational-function identities followed by the displayed tail asymptotics.

% END sections/06_worked_example.tex

% BEGIN sections/07_reproducibility.tex
\section{Reproducibility, proof dependencies, and verification}
\label{repro:section}

\subsection{What is supplied}

The package contains the modular article source, an assembled
single-file \TeX{} source, the compiled PDF, the exact
finite-field certificate, an explicit $(14,20)$ polynomial
identity with its finite telescoper, and portable verification
programs. A claim ledger separates proved results, inherited
inputs, numerical diagnostics, and remaining conjectures.
The source manifest pins every archived dependency to the
repository revision and a SHA-256 hash.

The main analytic dependency for gap six is the three-dimensional
Gamma basis and its complementary-minor normalization in the
ordered-boundary contribution \cite{OrderedIncoming}.
The recurrence and normalization are restated in the main proof;
Appendix~\ref{ga:appendix} proves the specialized Gamma contour
and initial-determinant assertions. The unchanged inherited source
is included for inspection under
\texttt{provenance/inherited/ordered\_gamma\_basis.tex}.
The earlier complementary-minor source
\cite{GapFourSource} and the two source sections posing the
complex-weight and independent-shift questions are also preserved.
They are provenance artifacts, and are not silently incorporated
as new authorship.

The new algebraic primitive theorem and its degree consequence
are proved directly in Section~\ref{cr:section}. The Hahn
transformation and Euler reductions are derived at the needed
normalization in Section~\ref{hp:section}. The entire unequal-rate
asymptotic is obtained by finite local subtraction in
Section~\ref{urc:section}; its coincident constants are attributed
to the inherited product formula and independently checked.

\subsection{Exact computations and their roles}

\begin{longtable}{@{}>{\raggedright\arraybackslash}p{0.31\textwidth}
 >{\raggedright\arraybackslash}p{0.61\textwidth}@{}}
\toprule
Program & What it establishes or checks\\
\midrule
\endfirsthead
\toprule
Program & What it establishes or checks\\
\midrule
\endhead
\texttt{verify\_gap\_six.py} &
The finite part of the proof: 2,044 nonzero determinants modulo
$2^{31}-1$, with trial-division primality and invertible
denominators. Also 2,049 independent cotangent coefficients
and 17 exact original-matrix normalization checks.\\[7pt]
{\small\texttt{create\_gap\_six\_example.py}} &
Generates the integer coefficient row, sextic subtraction,
and complete finite telescoper for $(p,q)=(14,20)$.
Checks exact rational differences and the limiting constant.\\[7pt]
\texttt{verify\_complex\_rank.py} &
66 exact assertions: rank-one minors, complex isotropic
updates, rational harmonic decompositions, exceptional cutoff
identities, and a varying-kernel pencil.\\[7pt]
\texttt{verify\_hahn\_polylogs.py} &
149 exact assertions: monic Hahn orthogonality and norms,
highest coefficients, compressed determinants, ensemble sums,
pole residues, and generator coefficients. Separate
75-digit quadratures test the generator and zeta endpoint.\\[7pt]
\texttt{derive\_euler\_endpoint.py} &
21 exact coefficient assertions derived from the height-one
gamma quotient and convergent shuffle/stuffle relations.
No numerical relation search is used.\\[7pt]
\texttt{verify\_collisions.py} &
119 exact assertions about primitives and contact coefficients.
Separate 65-digit beta and polygamma quadratures check the
complete collision subtraction at several rates.\\[7pt]
Two independent audit scripts &
Reconstruct original Bernoulli matrices, check the cofactor
normalization and exact threshold constants, and derive the
sensitive Stieltjes and trigamma finite parts independently.\\
\bottomrule
\end{longtable}

The modular certificate is finite and exhaustive on the stated
interval, because the analytic proof has already reduced the
unresolved indices to that interval. It would not prove an
unbounded conjecture without that analytic step. Conversely,
the high-precision integrations are not used to assert that an
unproved residual vanishes.

The independent audits rederived the complex rank-one minor,
all high-risk constants in the gap-six argument, the leading
inverse-row cofactor sign, and the locally uniform Stieltjes
remainder. In particular they recovered
$F_{1,0}^{0,0}=3\gamma_2/2+2\gamma\zeta(2)-\zeta(3)$ and
$F_{0,0}^{1,1}=2\zeta(3)+4\zeta'(3)$ separately from the
collision derivation. Audit notes and their exact output are
included. This is independent internal checking, not external
peer review or a formal proof-assistant certificate.

\subsection{Rebuild instructions}

The build is local and requires no network calls. In the package
directory run
\begin{verbatim}
python -m pip install -r requirements.txt
python verification/run_all.py
python build.py
\end{verbatim}
The verification dependencies are Python, SymPy, and mpmath.
The PDF build uses a standard \texttt{pdflatex} installation with
the packages listed in \texttt{article.tex}. The build program
assembles the standalone \TeX{} file and runs three passes for
the table of contents and references. It does not execute shell
escape.

The verification runner records each exit status and elapsed
time and stops on a failure. The result JSON files record their
scope and arithmetic mode. Numerical values can vary in their
last diagnostic digits with library versions; the exact
finite-field residue list is identified by its canonical hash.
The package manifest checks artifact integrity, while the
recurrences and the supplied programs make the mathematical
certificate independently reproducible.


% END sections/07_reproducibility.tex

% BEGIN references.tex
\begin{thebibliography}{99}
\raggedright
\addcontentsline{toc}{section}{References}

\bibitem{RepoManuscript}
ProveIt project, \emph{Polylogarithm manuscript}.
Repository revision \texttt{c4036d60715916ae72dd18983782fdb182840b86}.
\href{https://github.com/VladimirReshetnikov/ProveIt/tree/c4036d60715916ae72dd18983782fdb182840b86/Analysis/Polylogarithms/docs/manuscript}{Pinned manuscript directory}.
Inspected 11--12 October 2026.

\bibitem{RepoIncoming}
ProveIt project, \emph{Incoming research contributions}, same revision.
\href{https://github.com/VladimirReshetnikov/ProveIt/tree/c4036d60715916ae72dd18983782fdb182840b86/docs/incoming}{Pinned incoming directory}.
The full eleven-archive inventory and hashes are in the accompanying
provenance record.

\bibitem{OrderedIncoming}
ProveIt research continuation,
\emph{Ordered Boundary Identities and Gamma Kernels}, 11 October 2026.
Especially Section~8, \emph{The Gamma basis and eventual fixed-gap
reconstruction}, and Section~9, research status and questions.
\href{https://github.com/VladimirReshetnikov/ProveIt/blob/c4036d60715916ae72dd18983782fdb182840b86/docs/incoming/ProveIt_Ordered_Boundary_Identities_2026-10-11.zip}{Pinned archive}.
The Gamma-basis source used in this article is included unchanged
under \texttt{provenance/inherited/}.

\bibitem{BilinearIncoming}
ProveIt research continuation,
\emph{Bilinear Alternating Tails and Harmonic Determinants},
11 October 2026. Especially
\texttt{sections/05\_degree\_rank.tex} and
\texttt{sections/10\_research\_questions.tex}.
\href{https://github.com/VladimirReshetnikov/ProveIt/blob/c4036d60715916ae72dd18983782fdb182840b86/docs/incoming/ProveIt_Bilinear_Tails_and_Determinants_2026-10-11.zip}{Pinned archive}.

\bibitem{TailKernelsIncoming}
ProveIt research continuation,
\emph{Bilinear Tail Identities, Polylogarithmic Resonances,
and Stieltjes Boundary Products}, 11 October 2026.
\href{https://github.com/VladimirReshetnikov/ProveIt/blob/c4036d60715916ae72dd18983782fdb182840b86/docs/incoming/ProveIt_Tail_Kernels_and_Boundary_Products_2026-10-11.zip}{Pinned archive}.
This contribution already includes the smaller-order-twelve result.

\bibitem{LatticeIncoming}
ProveIt research continuation,
\emph{Lattice Contacts and Mixed Zeta Identities},
11 October 2026. Especially Sections~4 and~6, and Question~Q9
in \texttt{sections/08\_audit\_and\_research.tex}.
\href{https://github.com/VladimirReshetnikov/ProveIt/blob/c4036d60715916ae72dd18983782fdb182840b86/docs/incoming/ProveIt_Lattice_Contacts_and_Mixed_Zeta_Identities_2026-10-11.zip}{Pinned archive}.

\bibitem{GapFourSource}
ProveIt research continuation,
\emph{Harmonic Tail Gamma}:
\texttt{sections/02\_gap\_four.tex}.
\href{https://github.com/VladimirReshetnikov/ProveIt/blob/c4036d60715916ae72dd18983782fdb182840b86/Analysis/Polylogarithms/docs/reports/stieltjes-harmonic-resonance/continuations/harmonic-tail-gamma/sections/02_gap_four.tex}{Pinned recurrence and complementary-minor source}.

\bibitem{DLMF}
F.~W.~J. Olver and others, eds.,
\emph{NIST Digital Library of Mathematical Functions}.
Sections 5.11(ii) (complex Stirling remainder), 5.15 (polygamma),
18.2 (orthogonal polynomial zeros), 18.19--18.20 (Hahn
polynomials), and 25.11 (Hurwitz zeta).
\url{https://dlmf.nist.gov/}.
Accessed 11--12 October 2026.

\bibitem{ChenSinger2012}
S. Chen and M.~F. Singer,
\emph{Residues and telescopers for bivariate rational functions},
Advances in Applied Mathematics \textbf{49} (2012), 111--133.
\url{https://doi.org/10.1016/j.aam.2012.04.003}.
Preprint: \url{https://arxiv.org/abs/1203.4200}.
Proposition~2.5 is the rational discrete-residue criterion.

\bibitem{Krattenthaler2023}
C. Krattenthaler,
\emph{Hankel determinants of linear combinations of moments of
orthogonal polynomials, II},
The Ramanujan Journal \textbf{61} (2023), 597--627.
\url{https://doi.org/10.1007/s11139-021-00514-8}.
\url{https://arxiv.org/abs/2101.04225}.

\bibitem{Xu2017}
C. Xu, \emph{Multiple zeta values and Euler sums}.
Revised preprint, 2017.
\url{https://arxiv.org/abs/1609.05863}.
Related journal DOI: \url{https://doi.org/10.1016/j.jnt.2017.01.018}.

\bibitem{XuWang2019}
C. Xu and W. Wang,
\emph{Explicit formulas of Euler sums via multiple zeta values}.
Revised preprint, 2019.
\url{https://arxiv.org/abs/1805.08056}.

\bibitem{ShpotParis2016}
M.~A. Shpot and R.~B. Paris,
\emph{Integrals of products of Hurwitz zeta functions via Feynman
parametrization and two double sums of Riemann zeta functions}.
Revised preprint, 2020.
\url{https://arxiv.org/abs/1609.05658}.

\end{thebibliography}

% END references.tex

\end{document}
